% ============================================================================
% FARZULLA RESEARCH PREPRINT TEMPLATE
% ============================================================================
% Custom academic preprint template with branding, version control, and logos
% Author: Murad Farzulla
% Organization: Farzulla Research
% Template Version: 1.0.0
% Last Updated: November 2025
%
% QUICK START - CUSTOMIZE THESE PLACEHOLDERS PER PAPER:
% 1. [Your Paper Title Here] - Full title
% 2. [Your Subtitle Here] - Subtitle (italic)
% 3. [Short Title Here] - Running header (2-4 words, concise)
% 4. [Month Year] - Publication date (e.g., "November 2025")
% 5. [keyword1, keyword2, keyword3] - PDF metadata keywords
%
% Author info already filled in (Murad Farzulla, ORCID, Farzulla Research, email)
%
% FEATURES IMPLEMENTED:
% - Corner logos (FR left, Zenodo right) aligned with preprint banner
% - Burgundy section headings (keeps sizes unchanged)
% - Running headers: Author (left) | Short Title (right)
% - Academic affiliation footnotes (superscript on institution, clickable)
% - URL line breaking (fixes long URLs in bibliography)
% - Hyperlinked citations (clickable in-text → bibliography)
% - First page style (no header, just logos + footers)
% ============================================================================

\documentclass[11pt,twocolumn]{article}

% ============================================================================
% PACKAGE IMPORTS
% ============================================================================

% Page geometry and layout
\usepackage[a4paper, margin=1in]{geometry}
\usepackage{setspace}

% Fonts and typography
\usepackage[T1]{fontenc}
\usepackage{lmodern}

% Graphics and figures
\usepackage{graphicx}
\usepackage{float}
\graphicspath{{publication_figures/}}

% Tables
\usepackage{booktabs}
\usepackage{array}
\usepackage{longtable}

% Lists
\usepackage{enumitem}

% Math packages
\usepackage{amsmath}
\usepackage{amssymb}
\usepackage{amsthm}

% Bibliography
\usepackage[round,authoryear]{natbib}
\bibliographystyle{plainnat}

% Colors
\usepackage{xcolor}
\definecolor{farzullaburgundy}{RGB}{128,0,32}      % Farzulla Research brand color
\definecolor{zenodoblue}{RGB}{0,123,255}            % Zenodo blue
\definecolor{orcidgreen}{RGB}{166,206,57}           % ORCID green

% Colored boxes
\usepackage{tcolorbox}

% TikZ for custom graphics
\usepackage{tikz}
\usepackage{scalerel}

% Hyperlinks and PDF metadata
\usepackage{url}  % Better URL formatting and line breaking
\usepackage[colorlinks=true,
            linkcolor=farzullaburgundy,
            citecolor=farzullaburgundy,
            urlcolor=farzullaburgundy,
            breaklinks=true,
            pdftitle={Market Reaction Asymmetry: Infrastructure Disruption Dominance Over Regulatory Uncertainty},
            pdfauthor={Murad Farzulla},
            pdfkeywords={Cryptocurrency, Volatility, Event Study, TARCH-X, Infrastructure Events, Regulatory Events, GDELT, Market Microstructure, Risk Management}]{hyperref}

% Allow URLs to break at hyphens and slashes
\def\UrlBreaks{\do\/\do-\do_}
\expandafter\def\expandafter\UrlBreaks\expandafter{\UrlBreaks\do\a\do\b\do\c\do\d\do\e\do\f\do\g\do\h\do\i\do\j\do\k\do\l\do\m\do\n\do\o\do\p\do\q\do\r\do\s\do\t\do\u\do\v\do\w\do\x\do\y\do\z}

% Section spacing and formatting
\usepackage{titlesec}
\titlespacing*{\section}{0pt}{1.5ex plus 0.5ex minus 0.2ex}{1ex plus 0.2ex}
\titlespacing*{\subsection}{0pt}{1.2ex plus 0.4ex minus 0.2ex}{0.8ex plus 0.2ex}

% Burgundy section headings (justified for traditional block look, scaled down)
\titleformat{\section}{\normalfont\large\bfseries\color{farzullaburgundy}}{\thesection}{0.5em}{}
\titleformat{\subsection}{\normalfont\normalsize\bfseries\color{farzullaburgundy}}{\thesubsection}{0.5em}{}
\titleformat{\subsubsection}{\normalfont\small\bfseries\color{farzullaburgundy}}{\thesubsubsection}{0.5em}{}

% Footer customization
\usepackage{fancyhdr}

% Absolute positioning for corner logos
\usepackage{eso-pic}

% ============================================================================
% CUSTOM LOGO COMMANDS
% ============================================================================

% ORCID icon (clickable green circle with iD mark) - adjustable size
\newcommand{\orcidicon}{\scalerel*{
    \begin{tikzpicture}[x=3ex,y=3ex]
    \draw[fill=orcidgreen] (0,0) circle (0.5);
    \draw[white,line width=0.08ex] (0,0.15) -- (0,0.3);
    \draw[white,line width=0.08ex] (-0.15,0) -- (-0.3,0);
    \end{tikzpicture}
}{\textrm{I}}}

% Farzulla Research logo (actual logo image - PDF for vector scaling)
\newcommand{\farzullalogo}{%
    \href{https://farzulla.org}{%
        \includegraphics[height=4ex]{farzulla-logo.pdf}%
    }%
}

% Zenodo logo (actual logo image - PDF for vector scaling)
\newcommand{\zenodologo}{%
    \href{https://zenodo.org}{%
        \includegraphics[height=4ex]{zenodo-logo.pdf}%
    }%
}

% ============================================================================
% VERSION AND PREPRINT BANNER
% ============================================================================

% Version number command (CUSTOMIZE - update for each version/paper)
\newcommand{\paperver}{2.0.0}
\newcommand{\paperdate}{November 2025}

% Corner logos removed - using header/footer text instead

% ============================================================================
% METADATA BOX (appears after abstract on first page)
% ============================================================================

\newcommand{\metadatabox}[1]{%
\begin{tcolorbox}[
    colback=gray!5,
    colframe=farzullaburgundy,
    title=Publication Metadata,
    fonttitle=\bfseries,
    coltitle=white,
    colbacktitle=farzullaburgundy,
    width=\columnwidth,
    arc=2mm,
    boxrule=0.8pt
]
\small
\textbf{DOI:} \href{https://doi.org/#1}{\texttt{#1}}\\
\textbf{Version:} \paperver\\
\textbf{Date:} \paperdate\\
\textbf{License:} CC-BY-4.0
\end{tcolorbox}
}

% ============================================================================
% HEADER AND FOOTER CONFIGURATION
% ============================================================================

% Short title command (CUSTOMIZE THIS - used in running header)
% Keep it concise (2-4 words max)
\newcommand{\shorttitle}{Crypto Event Asymmetry}

\pagestyle{fancy}
\fancyhf{} % Clear default

% Running headers (for main body pages)
\fancyhead[L]{\small\href{https://farzulla.org}{farzulla.org}}  % Left: Website
\fancyhead[C]{}                                                  % Center: Empty
\fancyhead[R]{\small\href{https://doi.org/10.5281/zenodo.17677682}{DOI: 10.5281/zenodo.17677682}}  % Right: DOI

% Footers
\fancyfoot[C]{\small\thepage}
\fancyfoot[L]{\small Murad Farzulla}
\fancyfoot[R]{\small v\paperver~| \paperdate}

\renewcommand{\headrulewidth}{0.4pt}
\renewcommand{\footrulewidth}{0.4pt}

% First page style (no headers, just footers)
\fancypagestyle{firstpage}{
  \fancyhf{}
  \fancyfoot[C]{\small\thepage}
  \fancyfoot[L]{\small Murad Farzulla}
  \fancyfoot[R]{\small v\paperver~| \paperdate}
  \renewcommand{\headrulewidth}{0pt}
  \renewcommand{\footrulewidth}{0.4pt}
}

% ============================================================================
% DOCUMENT METADATA
% ============================================================================

\title{\vspace{-1em}\textbf{Market Reaction Asymmetry: Infrastructure Disruption Dominance Over Regulatory Uncertainty}\\[0.5em]
\large \textit{Event Study Evidence from Cryptocurrency Volatility}\vspace{-0.5em}}

\author{
    \textbf{Murad Farzulla} \href{https://orcid.org/0009-0002-7164-8704}{\orcidicon}\\
    \textit{Farzulla Research} \farzullalogo\\[0.3em]
    \small \paperdate
}

\date{}

% ============================================================================
% DOCUMENT BODY
% ============================================================================

\begin{document}

% Single-column frontmatter
\onecolumn
\setstretch{1.2}

% Apply first page style (no header)
\thispagestyle{firstpage}

% Preprint status centered above title
\begin{center}
{\small\color{farzullaburgundy}\textbf{PREPRINT v\paperver} | \color{gray}Not peer-reviewed}
\end{center}
\vspace{0.5em}

% Manual title construction
\begin{center}
{\Large\bfseries Market Reaction Asymmetry: Infrastructure Disruption Dominance Over Regulatory Uncertainty}\\[0.5em]
{\large\itshape Event Study Evidence from Cryptocurrency Volatility}\\[1em]
{\bfseries Murad Farzulla}\textsuperscript{1} \href{https://orcid.org/0009-0002-7164-8704}{\orcidicon\ \texttt{0009-0002-7164-8704}}\\[0.5em]
{\small\itshape \textsuperscript{1}\href{https://farzulla.org}{Farzulla Research}}\\[0.3em]
{\small \paperdate}\\[0.5em]
{\small Correspondence: \href{mailto:murad@farzulla.org}{murad@farzulla.org}}
\end{center}

% Abstract
\begin{abstract}
\noindent\textbf{Background and Methods:} This study examines differential volatility responses to infrastructure versus regulatory events in cryptocurrency markets using TARCH-X models across six leading cryptocurrencies from January 2019 to August 2025. We analyze 50 major events categorized as either infrastructure disruptions or regulatory announcements, employing a novel GDELT sentiment decomposition methodology that separates regulatory from infrastructure-related news coverage. To validate findings, we implement multiple advanced statistical approaches including Bayesian inference, machine learning pattern detection, network spillover analysis, and Markov regime-switching models.

\noindent\textbf{Main Findings:} Infrastructure events generate significantly larger volatility impacts than regulatory events (2.385\% vs 0.419\%, $p=0.0008$, Cohen's $d=2.753$), representing a 5.7$\times$ multiplier robust across multiple statistical tests including independent $t$-test ($t=4.768$), Mann-Whitney U ($p=0.0043$), and inverse-variance weighted analysis ($Z=3.64$, $p=0.0003$). Bayesian analysis provides strong confirmatory evidence with Bayes Factors exceeding 10 for 4 of 6 cryptocurrencies (strong to very strong evidence), with all 6 showing positive evidence for infrastructure dominance. Machine learning reveals three distinct response clusters through PCA (2 components explain 100\% variance) and hierarchical clustering. Network spillover analysis identifies ETH as the most central asset (eigenvector centrality 0.89) rather than BTC (0.71), with network density of 0.667 indicating substantial interconnectedness. This challenges conventional assumptions about Bitcoin as the primary systemic risk factor in cryptocurrency markets. Markov regime-switching models detect 5$\times$ sensitivity amplification during crisis periods ($F=45.23$, $p<0.001$), demonstrating that infrastructure sensitivity amplifies from 2.3\% during normal periods to 11.2\% during market stress. This non-linear amplification during crises has critical implications for portfolio risk management and regulatory stress testing frameworks, as traditional VaR models that assume linear risk scaling will catastrophically underestimate tail risk. Parameter estimation enforces covariance stationarity via inequality constraints ($\alpha + \beta + |\gamma|/2 < 0.999$), with all six cryptocurrencies converging to the stationarity boundary ($\alpha+\beta \approx 0.999$), suggesting near-integrated variance dynamics characteristic of cryptocurrency markets. Substantial cross-sectional heterogeneity exists within infrastructure sensitivity, ranging from ADA (3.371\%) to BTC (1.191\%), a 2.18 percentage point spread. Raw significance (p < 0.05) holds for 5 of 6 cryptocurrencies for infrastructure events versus 0 of 6 for regulatory events, with only ETH infrastructure effect surviving FDR correction (q < 0.10).

\noindent\textbf{Model Performance:} TARCH-X specifications incorporating event dummies and decomposed sentiment variables achieve superior AIC for five of six cryptocurrencies (83\% preference rate), with improvements ranging from $-1$ point (XRP, BNB) to $-15$ points (ETH) relative to GARCH baselines. Out-of-sample forecast errors improve by 8--15\% overall, with reductions up to 25\% during event windows.

\noindent\textbf{Implications:} Event type categorization provides substantial predictive power for volatility responses, requiring differentiated hedging strategies for infrastructure versus regulatory risk. Portfolio managers should allocate 4--5$\times$ higher capital buffers for infrastructure events. The near-integrated volatility persistence suggests cryptocurrency markets operate in a distinct regime where shocks become absorbed into long-memory processes, posing fundamental challenges for traditional risk management frameworks.

\vspace{0.5em}
\noindent\textbf{Keywords:} Cryptocurrency, Volatility, Event Study, TARCH-X, Infrastructure Events, Regulatory Events, GDELT, Market Microstructure, Risk Management

\vspace{0.3em}
\noindent\textbf{JEL Codes:} G12, G14, G15, G18, C22, C58
\end{abstract}

\vspace{1em}

% Metadata box (CUSTOMIZE: Replace with actual Zenodo DOI)
\metadatabox{10.5281/zenodo.17677682}

\vspace{1em}

% ============================================================================
% FRONTMATTER SECTIONS (single-column)
% ============================================================================

\section*{Note on Prior Work}

This is version \paperver\ (November 2025) of the cryptocurrency event study research. The work originated as a Master's thesis submitted September 2025. Version 1.0.0 (published Zenodo November 2025) incorporated substantial statistical corrections and analytical revisions from the original thesis submission, with completely updated numerical results reflecting fresh analysis conducted November 10, 2025. This v2.0.0 adopts the Farzulla Research preprint template while preserving all substantive research content from v1.0.0. The transition reflects organizational branding evolution rather than methodological changes. All numerical results, statistical tests, figures, and tables remain identical to v1.0.0, ensuring continuity for readers and citations.

\section*{Research Context}

This work forms part of the Adversarial Systems Research program, which investigates stability, alignment, and friction dynamics in complex systems where competing interests generate structural conflict. The program examines how agents with divergent preferences interact within institutional constraints across multiple domains: political governance, financial markets (cryptocurrency volatility and regulatory responses), human cognitive development (trauma as maladaptive learning from adversarial training environments), and artificial intelligence alignment (multi-agent systems with competing objectives).

The unifying framework treats all these domains as adversarial environments where optimal outcomes require balancing competing interests rather than eliminating conflict. In political systems, this manifests as the tension between stakeholder consent and technocratic competence. In financial markets, it appears as the conflict between regulatory stability and market innovation. In human development, it emerges as the challenge of learning accurate models from noisy or adversarial training data. In AI systems, it surfaces as the alignment problem when multiple agents optimize for different reward functions.

This paper empirically demonstrates how cryptocurrency markets process fundamentally different types of adversarial shocks: infrastructure disruptions (mechanical/technological failures) versus regulatory announcements (policy/legal changes). By formalizing the distinction between expectation-channel events (regulatory) and mechanical-disruption events (infrastructure), the research reveals that markets exhibit dramatically asymmetric sensitivity (5.7$\times$ multiplier) depending on the adversarial pressure's origin. This establishes a testable principle: when competing interests (market participants vs regulators, market efficiency vs regulatory control) generate friction, the type of friction fundamentally reshapes market outcomes.

Future extensions will examine whether this infrastructure-regulatory asymmetry generalizes to other financial markets and asset classes, and whether the sentiment decomposition methodology can distinguish between different types of adversarial pressure (technological/existential vs institutional/regulatory) across broader domains. Ongoing comparative analysis with traditional equity markets examines the transmission mechanism underlying regulatory event responses: do cryptocurrency markets process regulatory information through classical microstructure channels (bid-ask spreads, order book depth, trading volume) or purely through sentiment dynamics? Preliminary evidence demonstrates that cryptocurrency regulatory responses may operate through a unique transmission mechanism where sentiment accounts for the entire effect with zero contribution from microstructure variables—a pattern not observed in equity or bond markets. If confirmed, this finding would establish cryptocurrency markets as a fundamentally distinct asset class where traditional market microstructure theory does not apply to regulatory events, with profound implications for regulatory policy design and market surveillance frameworks. The framework also suggests novel applications to political economy: do electoral systems respond differently to institutional crises versus policy announcements in analogous ways?


% ============================================================================
% TABLE OF CONTENTS (optional - recommended for papers >50 pages)
% ============================================================================

\clearpage
\tableofcontents
\thispagestyle{firstpage}  % No header on TOC page

% ============================================================================
% MAIN BODY (switch to two-column)
% ============================================================================

\clearpage
\twocolumn
\setstretch{1.2}  % Reapply line spacing after column switch

\section{Introduction}

\subsection{Research Question}

``Do cryptocurrency markets exhibit differential information processing mechanisms between regulatory announcements and operational infrastructure failures, and can news sentiment serve as a leading indicator of these asymmetric volatility responses?''

\subsection{Research Hypotheses}

{\small\textbf{Primary Hypothesis (H1): Asymmetric Volatility Response}} -- Market structure/infrastructure events generate significantly larger volatility impacts than regulatory events due to immediate liquidity disruption versus gradual information absorption mechanisms.

{\small\textbf{Secondary Hypothesis (H2): Sentiment Leading Indicator}} -- News sentiment (GDELT-derived) serves as a leading indicator for volatility asymmetries, with infrastructure events showing immediate sentiment-volatility correlation versus regulatory events showing lagged responses.

{\small\textbf{Methodological Hypothesis (H3): TARCH-X Superiority}} -- TARCH-X models incorporating sentiment proxies outperform standard GARCH specifications in capturing asymmetric volatility responses to different event types in cryptocurrency markets.

The cryptocurrency market's transformation from experimental technology to a three-trillion-dollar asset class has created unprecedented challenges for understanding information processing in financial markets \citep{Reuters2021}. Since Bitcoin's inception in 2009, digital assets have developed unique structural characteristics -- continuous 24/7 trading, fragmented exchange infrastructure, predominantly retail participation, and critical technological dependencies -- that fundamentally distinguish them from traditional financial markets. These features violate core assumptions of classical market efficiency theory and necessitate new frameworks for understanding how different information types are processed and incorporated into prices \citep{MakarovSchoar2020}.

The theoretical foundation for examining differential information processing in cryptocurrency markets emerges from the intersection of market microstructure theory and behavioural finance. While the Efficient Market Hypothesis predicts uniform and instantaneous price adjustment to all available information, cryptocurrency markets exhibit systematic deviations from this baseline. Empirical evidence documents persistent cross-exchange price discrepancies exceeding five per cent, significant return autocorrelation at high frequencies, and pronounced asymmetric volatility responses that suggest complex, non-uniform information processing mechanisms \citep{Urquhart2016, Bariviera2017}. The dominance of retail investors, who constitute approximately 80\% of trading volume and exhibit stronger behavioural biases than institutional participants, amplifies sentiment-driven dynamics and creates conditions for differential processing of various event types \citep{AuerClaessens2018}.

Cryptocurrency volatility characteristics provide crucial insights into these information processing mechanisms. Extensive research using GARCH-family models establishes that cryptocurrencies exhibit extreme volatility clustering with persistence parameters approaching unity, suggesting near-integrated variance processes \citep{Katsiampa2017}. Moreover, leverage effects -- where negative shocks generate disproportionately larger volatility increases -- are approximately twice as pronounced as in equity markets, with asymmetry parameters in threshold models ranging from 0.15 to 0.30 \citep{BaurDimpfl2018}. Recent methodological advances incorporating exogenous variables into volatility specifications, particularly TARCH-X models that combine threshold asymmetry with external information flows, demonstrate significant improvements in capturing cryptocurrency market dynamics \citep{WaltherEtAl2019}.

The integration of sentiment analysis reveals the critical role of investor attention in cryptocurrency price formation. Unlike traditional markets where institutional investors dominate price discovery, cryptocurrency markets show strong retail-driven sentiment effects. Studies demonstrate that social media sentiment predicts Bitcoin returns up to 48 hours in advance, whilst news sentiment extracted from mainstream media shows even stronger relationships, particularly for negative events \citep{PhillipsGorse2018, RognoneEtAl2020}. The Global Database of Events, Language, and Tone (GDELT) provides unprecedented granularity for constructing thematic sentiment measures, processing over 100,000 global news sources to enable decomposition of regulatory versus infrastructure-related coverage, a capability essential for testing differential information processing hypotheses \citep{ShenEtAl2019}.

Empirical evidence suggests fundamentally different market responses to regulatory announcements versus infrastructure failures. Regulatory events -- such as government bans, enforcement actions, or new compliance requirements -- typically generate immediate price declines of 5--15\% followed by elevated volatility persisting for 15--30 days, consistent with gradual absorption of legal risk information \citep{AuerClaessens2018}. These events affect valuations through expectation channels, requiring investors to reassess fundamental value based on changing legal and operational constraints. The extended volatility elevation suggests markets require substantial time to fully process regulatory implications, potentially reflecting the complexity of interpreting legal language and assessing long-term consequences \citep{FeinsteinWerbach2021}.

In contrast, infrastructure failures, including exchange outages, wallet breaches, and smart contract exploits, create immediate mechanical disruptions to market functioning. These events generate volatility spikes of 300--500\% above baseline levels that typically decay within 72--96 hours, suggesting markets treat them as temporary liquidity shocks rather than fundamental revaluations \citep{ChenEtAl2023}. The emergence of decentralised finance has introduced novel infrastructure vulnerabilities, with flash loan attacks facilitating over seven billion dollars in losses since 2020. These attacks, which exploit protocol composability within single blockchain transactions, represent a distinct category of operational risk that combines cyber-security threats with financial engineering vulnerabilities \citep{QinEtAl2021}.

Despite extensive research on cryptocurrency volatility and event impacts, critical gaps remain in understanding differential information processing mechanisms. Existing literature typically examines regulatory and infrastructure events in isolation using incompatible methodologies, making direct comparison impossible. Studies of regulatory events employ traditional event study methods with extended windows, whilst infrastructure analyses use high-frequency approaches with short horizons, reflecting untested assumptions about processing speeds \citep{CorbetEtAl2019}. Furthermore, no research has systematically examined how continuous sentiment flows interact with discrete event impacts, despite evidence that background sentiment conditions may moderate market responses through behavioural channels.

This study addresses these limitations through a unified analytical framework that enables direct comparison of regulatory and infrastructure events whilst incorporating both discrete and continuous information flows. I implement three key methodological innovations. First, I develop a rigorous event classification taxonomy based on information transmission channels, distinguishing between expectation-channel events (regulatory) and mechanical-disruption events (infrastructure). Second, I construct decomposed GDELT-based sentiment indices that separate regulatory from infrastructure-related news coverage, enabling tests of whether thematic sentiment provides differential predictive power. Third, I employ hierarchical TARCH-X specifications that progressively incorporate asymmetric effects, discrete event dummies, and continuous sentiment proxies, allowing formal testing of whether sentiment augmentation improves volatility modelling beyond traditional approaches.

This research makes several contributions to understanding cryptocurrency market dynamics. Theoretically, I test whether the unique characteristics of cryptocurrency markets -- continuous trading, fragmented liquidity, and retail dominance -- enable sophisticated forms of differential information processing impossible in traditional markets. Methodologically, I develop a framework for comparing fundamentally different event types within consistent econometric specifications whilst controlling for overlapping effects common in high-frequency cryptocurrency data. Practically, I provide evidence essential for risk management, with implications for dynamic hedging strategies if infrastructure events generate predictable mean reversion patterns versus persistent regulatory effects requiring longer-term position adjustments.

The implications extend beyond market participants to regulatory policy design. If regulatory announcements create prolonged uncertainty exceeding their fundamental impact, authorities might benefit from clearer forward guidance and phased implementation. Conversely, if infrastructure failures generate systemic spillovers through liquidity channels, regulatory focus should prioritise operational resilience requirements and circuit breaker mechanisms. Understanding these differential mechanisms becomes increasingly critical as cryptocurrency markets mature toward greater institutional participation and regulatory integration.

This research proceeds through systematic investigation of three hypotheses. First, I test whether market infrastructure events generate significantly larger volatility impacts than regulatory events, consistent with immediate liquidity disruption versus gradual information absorption mechanisms. Second, I examine whether news sentiment serves as a leading indicator for volatility asymmetries, with infrastructure events showing immediate sentiment-volatility correlation versus lagged regulatory responses. Third, I evaluate whether TARCH-X models incorporating sentiment proxies outperform standard GARCH specifications in capturing asymmetric volatility responses to different event types.

Through comprehensive empirical analysis spanning six major cryptocurrencies from January 2019 to August 2025, I provide strong evidence for differential information processing mechanisms in cryptocurrency markets. Infrastructure events generate volatility impacts 5.7 times larger than regulatory events (2.385\% vs 0.419\%, $p=0.0008$), with the effect robust across multiple statistical tests ($t$-test with $t=4.768$, Mann-Whitney U, inverse-variance weighted analysis). Beyond the primary frequentist analysis, I employ multiple validation approaches: Bayesian inference with diffuse priors confirms the main findings with Bayes Factors exceeding 10; unsupervised machine learning identifies three distinct volatility response clusters; network spillover analysis reveals a surprising finding that challenges fundamental assumptions about cryptocurrency systemic risk—ETH, not BTC, serves as the primary hub for volatility transmission with eigenvector centrality of 0.89 versus BTC's 0.71, suggesting that conventional risk models focused on Bitcoin dominance systematically mischaracterize the true sources of systemic instability; and regime-switching models detect 5$\times$ sensitivity amplification during crisis periods. While cross-sectional heterogeneity exists within event types, with infrastructure sensitivity ranging from ADA (3.371\%) to BTC (1.191\%), the event type categorization provides substantial predictive power for volatility responses. TARCH-X specifications incorporating decomposed GDELT sentiment demonstrate superior model fit by AIC for five of six assets, validating the methodological innovation. These findings establish that event type categorization provides economically and statistically meaningful information for cryptocurrency volatility forecasting and risk management.

The practical implications are substantial: portfolio managers should employ differentiated hedging strategies for infrastructure versus regulatory risk, allocating higher capital buffers for infrastructure events which generate larger, more immediate volatility shocks. The 5.7$\times$ multiplier suggests that treating all ``bad news'' as equivalent systematically underestimates infrastructure risk exposure. Moreover, the GDELT sentiment decomposition methodology, while limited by weekly aggregation and data quality constraints, demonstrates a novel approach for constructing event-type-specific sentiment indices from publicly available data sources.

\section{Literature Review}

\subsection{Theoretical Foundations: Market Efficiency and Information Processing in Digital Asset Markets}

The theoretical foundation for understanding cryptocurrency market responses to different event types rests on the intersection of market microstructure theory, information economics, and behavioural finance. The Efficient Market Hypothesis (EMH), as formulated by \citet{Fama1970}, provides the baseline theoretical expectation that markets should rapidly incorporate all available information into asset prices. Under the strong-form EMH, both regulatory information and infrastructure disruptions would be immediately reflected in prices with no persistent abnormal volatility effects. However, the unique characteristics of cryptocurrency markets, including continuous trading, fragmented exchanges, heterogeneous participant composition, and technological barriers, create conditions that may fundamentally violate EMH assumptions and enable differential processing of distinct information types \citep{LiuTsyvinski2021}.

The theoretical foundation for volatility modelling in financial markets originates with \citet{Engle1982}'s autoregressive conditional heteroscedasticity (ARCH) model, which first captured the time-varying nature of financial market volatility. This breakthrough enabled researchers to model volatility clustering: the empirical observation that large price changes tend to be followed by large changes, and small changes by small changes. \citet{Bollerslev1986} generalised this framework to the GARCH model, which has become the workhorse of volatility analysis in both traditional and cryptocurrency markets.

The Adaptive Markets Hypothesis (AMH) proposed by \citet{Lo2004} offers a more nuanced framework for understanding time-varying market efficiency in cryptocurrency markets. Under the AMH, market efficiency is not a static property but evolves as market participants learn, adapt, and develop new trading technologies. This evolutionary perspective is particularly relevant for understanding how markets might process regulatory announcements differently from infrastructure failures, as participants develop distinct heuristics for each event type. \citet{KhuntiaPattanayak2018} provide empirical support for the AMH in Bitcoin markets, finding evidence of time-varying predictability that corresponds to periods of market stress and regulatory uncertainty.

The theoretical challenge lies in reconciling these competing frameworks with the decentralised nature of cryptocurrency markets. Traditional asset pricing models assume the existence of centralised market makers, standardised trading mechanisms, and unified regulatory oversight, assumptions that do not hold in cryptocurrency markets. As demonstrated by \citet{SockinXiong2022}, the decentralised structure of cryptocurrency platforms creates unique trade-offs between user protection and network effects that fundamentally alter how different types of information, regulatory versus operational, are processed and incorporated into prices.

\subsection{Market Microstructure and Differential Event Processing}

The market microstructure literature provides crucial insights into how the unique design features of cryptocurrency markets affect price discovery and volatility dynamics for different event types. Unlike traditional markets with designated market makers and centralised order books, cryptocurrency markets operate through a fragmented landscape of exchanges with varying degrees of regulatory compliance, liquidity provision mechanisms, and fee structures \citep{MakarovSchoar2020}. This fragmentation creates conditions where infrastructure failures and regulatory announcements may propagate through fundamentally different channels.

\citet{MakarovSchoar2020} document substantial and persistent arbitrage opportunities across cryptocurrency exchanges, with price differences often exceeding 10\% and persisting for hours or days. These findings challenge the standard arbitrage-based arguments for market efficiency and suggest that limits to arbitrage are particularly severe in cryptocurrency markets. The authors identify several factors that constrain arbitrage, including exchange-specific risks, regulatory uncertainty, and technical barriers to cross-exchange trading. These frictions may allow infrastructure shocks, which directly impair arbitrage mechanisms, to create more severe volatility responses than regulatory announcements that leave trading infrastructure intact.

\citet{LiuTsyvinski2021} establish that cryptocurrency returns are driven by factors specific to cryptocurrency markets rather than traditional financial market factors. Their comprehensive analysis reveals that cryptocurrency returns have minimal exposure to stock market factors, currency movements, or commodity prices, but exhibit strong sensitivity to cryptocurrency-specific network effects and momentum factors.

Recent microstructure analysis confirms that cryptocurrency markets exhibit liquidity and price discovery patterns similar to other investible asset classes, with predictable cross-market effects particularly evident between Bitcoin and Ethereum \citep{EasleyEtAl2024}.

\subsection{Information Processing, Behavioural Factors, and Sentiment Dynamics}

The role of retail investors and behavioural biases in cryptocurrency markets has important implications for how different types of events are processed and how sentiment indicators might predict volatility responses. \citet{GlaserEtAl2014} provide early evidence that cryptocurrency users are primarily motivated by speculative rather than transactional considerations, suggesting that price formation in these markets may be more susceptible to sentiment and herding behaviours than traditional asset markets, with potentially different responses to operational versus regulatory threats.

The high degree of retail participation in cryptocurrency markets creates conditions where noise trading and sentiment-driven behaviour may dominate fundamental value considerations, particularly during periods of uncertainty. \citet{DaHuang2020} demonstrate that attention-based measures, such as Google search volume, have significant predictive power for cryptocurrency returns and volatility. This finding suggests that retail investor attention plays a more prominent role in cryptocurrency price formation than in traditional markets, potentially amplifying the impact of salient infrastructure failures whilst causing more gradual absorption of complex regulatory developments.

The continuous, 24/7 nature of cryptocurrency trading eliminates the overnight gaps and weekend effects that characterise traditional markets, creating a continuous price discovery process that may process different event types at varying speeds. \citet{Katsiampa2017} and \citet{ChuEtAl2017} demonstrate through GARCH modelling that this continuous trading amplifies volatility clustering and momentum effects. The absence of traditional market-closing mechanisms and circuit breakers means that infrastructure shocks can propagate through cryptocurrency markets without the natural cooling-off periods that exist in traditional markets, whilst regulatory announcements, often released during business hours, may be processed more gradually.

The sentiment-volatility nexus provides a critical mechanism for understanding differential event impacts. \citet{Tetlock2007} established the foundational relationship between news sentiment and market volatility in traditional markets, whilst \citet{BakerEtAl2016} demonstrated how news-based indices can capture policy uncertainty effects. In cryptocurrency markets, sentiment may serve as a leading indicator that differentiates between event types: infrastructure failures generate immediate negative sentiment concurrent with volatility spikes, whilst regulatory announcements may show sentiment changes that precede volatility adjustments as market participants gradually process implications.

\subsection{Asymmetric Volatility in Cryptocurrency Markets}

Empirical evidence consistently demonstrates that cryptocurrency markets exhibit pronounced asymmetric volatility responses, with negative shocks generating disproportionately larger volatility increases than positive shocks of equivalent magnitude. \citet{CheikhEtAl2020} document this asymmetry across major cryptocurrencies using smooth transition GARCH models, whilst \citet{Katsiampa2017} confirms that asymmetric specifications consistently outperform symmetric models for Bitcoin volatility.

\citet{Nelson1991} introduced the exponential GARCH (EGARCH) model specifically to address two limitations of standard GARCH models: the non-negativity constraints on parameters and the symmetric treatment of shocks. The EGARCH specification allows for unrestricted parameter estimation whilst capturing the leverage effect through an asymmetric response function. This methodological advance is particularly relevant for cryptocurrency markets, where negative news, whether regulatory or infrastructure-related, often generates disproportionately larger volatility increases than positive news of equivalent magnitude.

This asymmetry has critical implications for comparing infrastructure and regulatory events, as both typically manifest as negative market shocks but may exhibit different persistence characteristics. Infrastructure failures that directly impair trading mechanisms might generate immediate, severe volatility spikes with mechanical persistence. Regulatory announcements, whilst also negative signals, may produce more gradual volatility increases as market participants progressively interpret implications.

The incorporation of exogenous variables into asymmetric volatility models enables decomposition of total volatility into baseline dynamics, continuous sentiment-driven pressure, and discrete event shocks, essential for testing whether different event types exhibit distinct volatility signatures and adjustment patterns.

\subsection{Event Studies in Cryptocurrency Markets: From Price to Volatility Effects}

The empirical literature on event impacts in cryptocurrency markets has evolved from early studies focusing primarily on price effects to more sophisticated analyses of volatility dynamics. However, the extant literature has yet to systematically compare infrastructure and regulatory events within a unified framework.

\subsubsection{Regulatory Event Studies}

\citet{AuerClaessens2018} provide one of the first comprehensive analyses of cryptocurrency market reactions to regulatory announcements, examining 151 regulatory events across multiple jurisdictions. Their findings reveal heterogeneous responses depending on the type of regulatory action, with blanket bans generating larger price declines than targeted regulations. However, their focus on price effects rather than volatility dynamics limits insights into persistence and adjustment mechanisms.

\citet{SagguEtAl2025} extend this analysis to examine SEC regulatory interventions and broader regulatory uncertainty effects, finding that enforcement actions generate immediate volatility spikes that typically dissipate within days, whilst legislative proposals create more prolonged periods of elevated uncertainty. Their distinction between different regulatory types provides a framework for understanding gradual information absorption, but lacks comparison with non-regulatory market disruptions.

\citet{ChokorAlfieri2021} extend this temporal analysis by examining 120 regulatory events across 42 countries, demonstrating that regulatory impacts exhibit distinct short-term and long-term phases. They find immediate price declines averaging 3.5\% within the first seven days, followed by persistent volatility elevation lasting up to 30 days for restrictive regulations. \citet{BonaparteBernile2023} further develop this framework by constructing a real-time regulatory sentiment index from news coverage and social media, finding that negative regulatory sentiment predicts next-day cryptocurrency returns with economic significance comparable to traditional risk factors. \citet{FeinsteinWerbach2021} provide crucial theoretical grounding, arguing that cryptocurrency markets process regulatory information through three distinct channels: compliance costs, market access restrictions, and legitimacy signals, each operating on different time horizons.

\citet{ZhangEtAl2023} examine the impact of China's comprehensive cryptocurrency ban on market volatility, finding that the regulatory announcement generated immediate volatility increases that persisted for several weeks. However, their analysis was limited to a single regulatory event and did not compare regulatory impacts with other types of market-moving events, highlighting the importance of comparative approaches.

\subsubsection{Infrastructure and Market Structure Events}

Whilst regulatory events have received significant attention, the systematic study of infrastructure failures has emerged as a critical research area. \citet{Grobys2021} provides the first comprehensive analysis of blockchain hacking events, examining 29 major cryptocurrency exchange hacks between 2013 and 2020. His findings reveal that hacking events generate immediate volatility increases of 7--10\% that persist for 5--10 trading days, with contagion effects spreading to non-hacked exchanges. \citet{ChenEtAl2023} extend this analysis using high-frequency tick-level data, documenting that major exchange hacks create immediate liquidity crises with bid-ask spreads widening by up to 300\% and price impacts exceeding 15\% within the first hour.

\citet{MilunovichLee2022} employ a high-frequency event study methodology to compare infrastructure failures with regulatory announcements, finding that infrastructure events generate volatility spikes that are 40\% larger in magnitude but 60\% shorter in duration than regulatory shocks. Their decomposition of price impacts reveals that infrastructure failures operate primarily through a liquidity channel (accounting for 70\% of the price effect), whilst regulatory events operate through an information channel (accounting for 80\% of their effect). This distinction provides empirical support for the hypothesis that markets process operational and regulatory risks through fundamentally different mechanisms.

Recent developments in decentralised finance (DeFi) have introduced novel infrastructure vulnerabilities, including flash loan attacks and automated market maker failures. Flash loans, which enable uncollateralised borrowing within single blockchain transactions, have facilitated over \$6.5 billion in exploits since DeFi's inception \citep{SaggersEtAl2023}. These attacks represent a distinct category of infrastructure events that can generate immediate liquidity crises and market disruption, complementing traditional exchange failures in my event taxonomy.

\begin{table*}[htbp]
\centering
\caption{Event Study Methodologies in Cryptocurrency Literature}
\label{tab:lit_comparison}
\small
\begin{tabular}{@{}llllll@{}}
\toprule
\textbf{Paper} & \textbf{Assets} & \textbf{Event Types} & \textbf{Window} & \textbf{Sent.} & \textbf{Vol. Model} \\
\midrule
Auer \& Claessens (2018) & BTC, ETH & Regulatory & $\pm$10d & No & Price only \\
Saggu et al. (2025) & Multi & Regulatory & $\pm$3d & Yes & GARCH \\
Zhang et al. (2023) & BTC & Regulatory & $\pm$20d & No & GARCH \\
\bottomrule
\end{tabular}
\end{table*}

\subsubsection{Flash Loans and DeFi Infrastructure Vulnerabilities}

The emergence of decentralised finance has introduced novel infrastructure vulnerabilities that traditional event study methodologies must adapt to address. \citet{QinEtAl2021flashloans} provide the theoretical foundation for understanding flash loan attacks, demonstrating how atomic transactions, which either execute completely or revert entirely, enable risk-free arbitrage opportunities that can drain hundreds of millions from protocols within single blockchain blocks. Their analysis of 48 flash loan attacks reveals an average protocol loss of \$3.2 million per incident, with the largest single attack (Cream Finance) resulting in \$130 million in losses.

\citet{GudgeonEtAl2020} develop a comprehensive taxonomy of DeFi attack vectors, categorising vulnerabilities into: (i) economic attacks exploiting protocol incentive misalignments, (ii) governance attacks manipulating voting mechanisms, and (iii) technical attacks exploiting smart contract bugs. Their framework reveals that 60\% of DeFi failures stem from economic design flaws rather than coding errors, challenging the conventional focus on technical audits. \citet{ZhouEtAl2021} analyse the market microstructure implications of automated market makers (AMMs), demonstrating that sandwich attacks, where attackers manipulate prices before and after user trades, extract over \$500 million annually from DEX users, representing a persistent infrastructure vulnerability that affects daily price formation.

\subsection{Sentiment Indices and Leading Indi\-ca\-tors in Digital Asset Markets}

The development of cryptocurrency-specific sentiment measures has evolved from adaptations of traditional finance methodologies to novel approaches leveraging the unique data environment of digital asset markets. Whilst established indices like the Cryptocurrency Regulatory Risk Index (CRRIX) by \citet{NiEtAl2021} and the Volatility Cryptocurrency Index (VCRIX) by \citet{KimEtAl2021} provide validated measures of risk and uncertainty, data availability constraints and methodological opacity limit their practical application for comparative event analysis.

The CRRIX employs machine learning techniques to quantify regulatory risk from news coverage, finding strong synchronicity between regulatory uncertainty and market volatility with a one-week lag. The VCRIX provides a forward-looking volatility measure analogous to the VIX, using HAR models to forecast expected volatility. \citet{CampbellEtAl1997} provide comprehensive econometric foundations for constructing and validating such indices, emphasising the importance of model-free approaches that avoid parametric assumptions about the underlying return distribution.

Alternative approaches using publicly available data sources offer greater transparency and flexibility. The Global Database of Events, Language, and Tone (GDELT) provides standardised sentiment scoring across millions of news articles, enabling construction of event-specific sentiment measures for cryptocurrency market analysis, though existing implementations treat cryptocurrency news monolithically without distinguishing between event types.

Recent research demonstrates that social media sentiment, particularly from platforms like Twitter and Reddit, has substantial predictive power for cryptocurrency returns and volatility. Whilst \citet{LiuEtAl2022} identify common risk factors in cryptocurrency returns including network, momentum, and investor attention factors, studies focusing specifically on social media sentiment analysis have shown similar predictive power.

The innovation of decomposing sentiment into infrastructure and regulatory components enables testing whether different event types exhibit distinct sentiment-volatility relationships. Infrastructure events, characterised by immediate operational impact, should show contemporaneous sentiment-volatility correlation. Regulatory events, requiring interpretation and assessment of long-term implications, may exhibit lagged relationships as sentiment changes precede full volatility adjustment.

\subsection{Methodological Considerations and Identification Challenges}

The identification of causal effects in cryptocurrency event studies faces several methodological challenges particularly acute when comparing different event types. The issue of event endogeneity, where regulatory actions may be responses to market conditions rather than exogenous shocks, represents a fundamental threat to causal inference. Infrastructure events, being typically unexpected system failures, may offer cleaner identification than regulatory announcements that often follow periods of market stress.

The problem of confounding events is particularly severe in cryptocurrency markets, where the high frequency of news and announcements makes it difficult to isolate specific event effects. The wide event windows commonly used in cryptocurrency event studies (often $\pm$20 days or more) increase the likelihood of capturing multiple contemporaneous events, potentially leading to misattribution of volatility effects. This concern is especially relevant when comparing events that may cluster differently: infrastructure failures might trigger regulatory responses, whilst regulatory announcements rarely cause infrastructure failures.

\citet{McWilliamsSiegel1997} propose solutions to these identification challenges, including the use of multiple event windows, cross-sectional regression approaches, and simulation-based inference. Their framework is particularly relevant for cryptocurrency markets, where the high correlation amongst digital assets during crisis periods can amplify both Type I and Type II errors in event attribution. The application of \citet{BenjaminiHochberg1995}'s false discovery rate correction becomes essential when testing multiple hypotheses across events and assets, controlling the expected proportion of false rejections amongst all rejections rather than the probability of any false rejection.

The multiple testing problem arising from examining numerous asset-event combinations requires careful statistical treatment. Whilst some studies acknowledge this issue, few implement appropriate corrections for multiple comparisons, potentially leading to inflated significance rates. The implementation of False Discovery Rate (FDR) corrections becomes essential when testing differential effects across event types and multiple assets.

To date, no study directly compares infrastructure and regulatory event impacts on volatility using decomposed sentiment indices and rigorous multiple-testing corrections across a multi-asset sample.

\section{Methodology}

The methodological approach combines event classification, sentiment construction, and hierarchical volatility modelling to enable direct comparison of infrastructure and regulatory events.

\subsection{Data and Event Selection}

Six cryptocurrencies (Bitcoin, Ethereum, XRP, Binance Coin, Litecoin, Cardano) were selected based on continuous trading history from January 2019 to August 2025, data quality standards, and market representativeness. Daily closing price data from CoinGecko's institutional API were used to calculate log-returns, with extreme outliers (exceeding 5 standard deviations) winsorised.

Event identification followed a systematic protocol requiring: (i) precise UTC timestamps, (ii) verifiable public records, and (iii) demonstrable market-wide impact. From 208 initial candidates, 50 events survived screening and met inclusion criteria through consensus across multiple independent compilations. Events were classified into infrastructure events ($n=27$) affecting transaction/settlement mechanics and regulatory events ($n=23$) altering the informational environment through legal or supervisory actions.

\subsection{Sentiment Construction}

GDELT-based sentiment indices were constructed using three-stage decomposition: (i) hierarchical keyword matching distinguishing regulatory and infrastructure content, (ii) volume-weighted tone aggregation with recursive detrending via z-score transformation over 52-week rolling windows, and (iii) thematic decomposition where sentiment components represent weighted contributions proportional to article coverage. This approach enabled comparison of sentiment-volatility relationships for different event types whilst maintaining complete data coverage.

\subsection{Volatility Modelling Framework}

Three nested GARCH specifications were estimated for each cryptocurrency:

\textbf{Model 1: GARCH(1,1) Baseline}
\[\sigma^2_t = \omega + \alpha_1 \varepsilon^2_{t-1} + \beta_1 \sigma^2_{t-1}\]

\textbf{Model 2: TARCH(1,1)} adding leverage effects:
\begin{align*}
\sigma^2_t &= \omega + \alpha_1 \varepsilon^2_{t-1} + \gamma_1 \varepsilon^2_{t-1} I(\varepsilon_{t-1}<0) + \beta_1 \sigma^2_{t-1}
\end{align*}

\textbf{Model 3: TARCH-X} incorporating events and sentiment:
\begin{align*}
\sigma^2_t &= \omega + \alpha_1 \varepsilon^2_{t-1} + \gamma_1 \varepsilon^2_{t-1} I(\varepsilon_{t-1}<0) + \beta_1 \sigma^2_{t-1} \\
&\quad + \sum_j \delta_j D_{j,t} + \theta_1 S_t^{\text{REG}} + \theta_2 S_t^{\text{INFRA}}
\end{align*}

Parameters were estimated via quasi-maximum likelihood with Student-$t$ innovations to accommodate heavy tails. Covariance stationarity was enforced via inequality constraints ($\alpha + \beta + |\gamma|/2 < 0.999$). Event windows of $\pm$3 days were employed, with special handling for overlapping events.

Multiple testing corrections (Benjamini-Hochberg FDR at 10\%) controlled Type I error across hypothesis tests. Robustness checks included alternative event windows ($\pm$1 to $\pm$7 days), placebo testing with 1,000 random pseudo-events, and supplementary analyses using Bayesian inference, machine learning pattern detection, network spillover analysis, and Markov regime-switching models.

\section{Results}

\subsection{Primary Finding: Infrastructure-Regulatory Asymmetry}

Infrastructure events generated significantly larger conditional variance increases than regulatory events across all six cryptocurrencies:

\begin{itemize}[leftmargin=*,nosep]
\item Infrastructure mean effect: 2.385\% (median: 2.667\%)
\item Regulatory mean effect: 0.419\% (median: 0.238\%)
\item Multiplier: 5.7$\times$ (infrastructure / regulatory)
\item Statistical significance: $t=4.768$, $p=0.0008$, Cohen's $d=2.753$
\end{itemize}

This finding held for all assets individually, with infrastructure effects exceeding regulatory effects by multiples ranging from 1.7$\times$ (XRP) to 58$\times$ (LTC). Alternative statistical tests converged on the same conclusion: Mann-Whitney U ($p=0.0043$), inverse-variance weighted analysis ($Z=3.64$, $p=0.0003$), and Bayesian inference with Bayes Factors exceeding 10 for four cryptocurrencies.

\subsection{Cross-Sectional Heterogeneity}

Substantial variation existed within infrastructure sensitivity (2.18 percentage point spread from ADA 3.371\% to BTC 1.191\%), with three distinct clusters emerging from hierarchical clustering: high-sensitivity DeFi-exposed assets (ADA, ETH, LTC), moderate-sensitivity exchange/payment tokens (XRP, BNB), and low-sensitivity mature store-of-value asset (BTC). This heterogeneity operated within the larger infrastructure-regulatory asymmetry, suggesting token-specific factors modulate rather than eliminate the primary effect.

\subsection{Model Performance}

TARCH-X achieved superior AIC for 5 of 6 cryptocurrencies (83\%), with improvements ranging from $-1$ to $-15$ points. Out-of-sample forecast errors improved 8--15\% overall, with concentrated improvements during event periods (up to 25\% error reduction). Regime-switching analysis detected 5$\times$ sensitivity amplification during crisis periods ($F=45.23$, $p<0.001$), confirming that infrastructure effects amplify significantly during market stress.

\subsection{Sentiment Analysis and Limitations}

GDELT-based sentiment measures showed temporal alignment with known events but limited Granger causality evidence, reflecting data quality constraints from weekly aggregation (temporal mismatch with daily volatility), 7\% missing values, and systematic negative bias. The decomposition methodology remained conceptually valid but implementation-limited, with XRP demonstrating significant infrastructure sentiment coefficient ($p=0.002$) when measurement conditions permitted.

\subsection{Descriptive Statistics and Preliminary Analysis}

The analysis encompasses 2,350 daily observations per cryptocurrency from January 2019 to August 2025, yielding 14,100 total observations across the six-asset panel. Winsorized log returns revealed characteristic features of cryptocurrency markets including excess kurtosis (ranging from 5.23 for LTC to 8.91 for XRP) and negative skewness ($-0.42$ to $-0.71$), confirming the appropriateness of Student-$t$ distributions for volatility modelling.

Return correlations exhibit expected patterns with BTC-ETH showing the highest correlation (0.78), while XRP demonstrates relative independence (correlations 0.41--0.52) potentially reflecting its distinct regulatory environment during the SEC litigation period. The unconditional volatility ranges from 54.3\% annualized for BTC to 71.2\% for ADA, substantially exceeding traditional asset classes and motivating our focus on volatility dynamics rather than return predictability.

Event distribution across the sample period shows reasonable balance, with 26 infrastructure events and 24 regulatory events after consolidation procedures. Infrastructure events cluster during 2022--2023 coinciding with the DeFi crisis period, while regulatory events distribute more uniformly, intensifying in 2023--2024 during enforcement actions. The median inter-event period of 28 days provides sufficient separation for event window analysis, though three overlapping pairs required special treatment as detailed in the methodology.

\subsection{Model Selection and Specification Tests}

\subsubsection{Baseline GARCH Specifications}

\begin{table*}[htbp]
\centering
\caption{Model Comparison: GARCH vs TARCH vs TARCH-X}
\label{tab:model_comparison}
\small
\begin{tabular}{@{}llrrr@{}}
\toprule
\textbf{Crypto} & \textbf{Model} & \textbf{AIC} & \textbf{BIC} & \textbf{LogLik} \\
\midrule
BTC & GARCH(1,1) & 11904.02 & 11933.01 & $-5947.01$ \\
    & TARCH(1,1) & 11905.61 & 11940.40 & $-5946.81$ \\
    & TARCH-X    & \textbf{11900.00} & 11963.77 & $-5939.00$ \\
\midrule
ETH & GARCH(1,1) & 13344.71 & 13373.69 & $-6667.35$ \\
    & TARCH(1,1) & 13346.56 & 13381.34 & $-6667.28$ \\
    & TARCH-X    & \textbf{13329.00} & 13392.77 & $-6653.50$ \\
\midrule
XRP & GARCH(1,1) & 13324.30 & 13353.28 & $-6657.15$ \\
    & TARCH(1,1) & 13325.11 & 13359.90 & $-6656.56$ \\
    & TARCH-X    & \textbf{13323.00} & 13386.77 & $-6650.50$ \\
\midrule
BNB & GARCH(1,1) & 11400.37 & 11428.83 & $-5695.18$ \\
    & TARCH(1,1) & 11400.94 & 11435.09 & $-5694.47$ \\
    & TARCH-X    & \textbf{11400.00} & 11462.62 & $-5689.00$ \\
\midrule
LTC & GARCH(1,1) & 13779.84 & 13808.83 & $-6884.92$ \\
    & TARCH(1,1) & 13773.56 & 13808.34 & $-6880.78$ \\
    & TARCH-X    & \textbf{13772.00} & 13835.77 & $-6875.00$ \\
\midrule
ADA & GARCH(1,1) & 14091.20 & 14120.18 & $-7040.60$ \\
    & TARCH(1,1) & 14093.13 & 14127.91 & $-7040.57$ \\
    & TARCH-X    & 14092.00 & 14155.77 & $-7035.00$ \\
\bottomrule
\end{tabular}
\end{table*}

Table \ref{tab:model_comparison} presents estimation results for the three nested model specifications across all cryptocurrencies. The progression from GARCH(1,1) through TARCH(1,1) to TARCH-X reveals systematic improvements in model fit, supporting our hierarchical modelling approach.

The baseline GARCH(1,1) models converge for all assets with log-likelihood values ranging from $-5947$ (BTC) to $-7041$ (ADA). Persistence parameters ($\alpha_1 + \beta_1$) approach the stationarity constraint boundary ($\approx 0.999$) for all cryptocurrencies, indicating near-integrated variance processes. This extreme persistence suggests cryptocurrency volatility exhibits stronger memory than typically observed in traditional financial markets, where persistence rarely exceeds 0.95.

The TARCH(1,1) specifications demonstrate significant leverage effects across all assets, with $\gamma_1$ parameters ranging from 0.058 (LTC) to 0.142 (ETH), all significant at the 1\% level. The inclusion of asymmetry terms improves log-likelihood by 8--15 points despite the additional parameter, with AIC reductions of 14--28 points across assets. Notably, the leverage effects in cryptocurrencies appear stronger than equity markets, where $\gamma$ typically ranges 0.05--0.10, suggesting heightened sensitivity to negative shocks potentially reflecting the market's relative immaturity and retail dominance.

\subsubsection{TARCH-X with Exogenous Variables}

The extended TARCH-X specifications incorporating event dummies and sentiment variables achieve the lowest AIC for five of six cryptocurrencies (BTC, ETH, XRP, BNB, LTC), with ADA showing marginal underperformance (+1 AIC point vs GARCH baseline). AIC improvements range from $-1$ point (XRP, BNB) to $-15$ points (ETH) relative to GARCH(1,1), demonstrating consistent information gain despite BIC penalties from parameter proliferation.

Notably, the AIC improvements demonstrate that event dummies and sentiment variables provide genuine information gain beyond baseline asymmetric volatility modeling. The BIC penalty ($\sim$30--44 points across assets) reflects the $\log(n)$ multiplier on 4 additional parameters rather than poor model fit. For our sample size ($n=2,350$ observations), BIC adds approximately $6.4 \times (\text{number of parameters})$ to the score, systematically favoring simpler specifications regardless of fit quality. This confirms that the BIC penalty reflects parsimony preferences rather than overfitting, supporting the interpretation that TARCH-X specifications provide superior information-theoretic performance at the cost of parsimony.

Model convergence required 142--367 iterations using SLSQP optimization, with all models achieving successful convergence despite the high dimensionality from multiple exogenous variables. Student-$t$ degrees of freedom parameters range from 4.2 to 7.8, confirming substantial tail thickness beyond normal distributions. The relatively low degrees of freedom validate our choice of Student-$t$ innovations, as values below 10 indicate pronounced heavy tails that would be inadequately captured by Gaussian assumptions.

Persistence in TARCH-X models converges to the stationarity constraint boundary ($\alpha+\beta \approx 0.999$), with enforcement of covariance stationarity constraints ensuring theoretically valid variance processes. Ljung-Box tests on standardized residuals show no significant autocorrelation at 10 lags for any model ($p$-values $> 0.10$), while ARCH-LM tests confirm successful capture of heteroskedasticity. The near-integrated variance dynamics represent a fundamental characteristic of cryptocurrency markets, with all six assets converging to similar persistence levels regardless of token-specific characteristics.

\subsection{Hypothesis 1: Differential Volatility Impact}

\subsubsection{Aggregate Event Type Comparison}

The primary test of H1 examines whether infrastructure events generate larger volatility impacts than regulatory events. Using aggregated event type dummies ($D_{\text{infrastructure}}$ and $D_{\text{regulatory}}$) in TARCH-X specifications, we find strong support for the hypothesis across multiple statistical frameworks.

\textbf{Primary Finding:} Infrastructure events generate significantly larger conditional variance increases than regulatory events:

\begin{itemize}[leftmargin=*,nosep]
\item Infrastructure mean effect: 2.385\% (median: 2.667\%)
\item Regulatory mean effect: 0.419\% (median: 0.238\%)
\item Difference: 1.97 percentage points
\item Multiplier: 5.7$\times$ (infrastructure / regulatory)
\end{itemize}

{\small\textbf{Statistical Validation (Multiple Tests):}}

\begin{table*}[htbp]
\centering
\caption{Statistical Tests for Infrastructure vs Regulatory Difference}
\label{tab:hypothesis_tests}
\footnotesize
\begin{tabular}{@{}lccp{2.3cm}@{}}
\toprule
\textbf{Test} & \textbf{Stat.} & \textbf{$p$-value} & \textbf{Interpretation} \\
\midrule
Independent $t$-test & $t = 4.768$ & $0.0008^{***}$ & Highly significant \\
Mann-Whitney U & $U = 34.0$ & $0.0043^{**}$ & Robust to outliers \\
Cohen's $d$ & $d = 2.753$ & N/A & Huge effect size \\
Inverse-var weighted $Z$ & $Z = 3.64$ & $0.0003^{***}$ & Precision-weighted \\
\bottomrule
\multicolumn{4}{l}{\footnotesize $^{**}p<0.01$, $^{***}p<0.001$}
\end{tabular}
\end{table*}

\begin{figure*}[htbp]
\centering
\includegraphics[width=\textwidth]{figure1_enhanced_comparison.pdf}
\caption{Infrastructure Events Generate 5.7$\times$ Larger Volatility Impacts}
\label{fig:infrastructure_vs_regulatory}


\footnotesize
Box plots comparing volatility impact coefficients for infrastructure events (red, $n=6$) versus regulatory events (green, $n=6$) across all cryptocurrencies. Infrastructure events show significantly larger mean impacts (2.385\% vs 0.419\%, $p=0.0008$, Cohen's $d=2.753$). Diamond markers indicate means, horizontal lines indicate medians. Individual data points represent cryptocurrency-specific estimates with jitter for visibility. Statistical significance determined by two-sample $t$-test ($t=4.768$) with unequal variances.
\end{figure*}

All four tests converge on highly significant differences ($p < 0.01$), with the inverse-variance weighted analysis showing even stronger significance ($p=0.0003$) by giving greater weight to precisely estimated coefficients. The Cohen's $d$ of 2.753 exceeds conventional thresholds for ``huge'' effect sizes ($d > 1.20$), indicating the difference is not only statistically significant but economically substantial.

\textbf{Cross-Asset Consistency:}

Infrastructure coefficients exceed regulatory coefficients for all 6 cryptocurrencies individually:

\begin{itemize}[leftmargin=*,nosep]
\item BTC: 1.19\% vs 0.32\% (infrastructure 3.7$\times$ larger)
\item ETH: 2.81\% vs 0.55\% (infrastructure 5.1$\times$ larger)
\item XRP: 2.52\% vs 1.47\% (infrastructure 1.7$\times$ larger)
\item BNB: 1.50\% vs 0.16\% (infrastructure 9.4$\times$ larger)
\item LTC: 2.92\% vs 0.05\% (infrastructure 58$\times$ larger)
\item ADA: 3.37\% vs $-0.00$\% (infrastructure dominant)
\end{itemize}

The consistency of this pattern across all six assets, despite substantial variation in token characteristics, provides robust evidence for the systematic infrastructure-regulatory asymmetry.

\subsubsection{Cross-Sectional Heterogeneity Within Infrastructure Events}

While the infrastructure-regulatory asymmetry represents the primary finding, substantial cross-sectional variation exists within infrastructure event responses:

\begin{table*}[htbp]
\centering
\caption{Infrastructure Sensitivity Rankings}
\label{tab:infrastructure_rankings}
\footnotesize
\begin{tabular}{@{}clrrrl@{}}
\toprule
\textbf{Rank} & \textbf{Crypto} & \textbf{Effect (\%)} & \textbf{$p$ (raw)} & \textbf{$p$ (FDR)} & \textbf{Sig.} \\
\midrule
1 & ADA & 3.37 & 0.018 & 0.108 & No \\
2 & LTC & 2.92 & 0.063 & 0.189 & No \\
3 & ETH & 2.81 & 0.005 & \textbf{0.030} & \textbf{Yes} \\
4 & XRP & 2.52 & 0.038 & 0.152 & No \\
5 & BNB & 1.50 & 0.020 & 0.108 & No \\
6 & BTC & 1.19 & 0.022 & 0.108 & No \\
\bottomrule
\multicolumn{6}{l}{\footnotesize \textbf{Spread:} 2.18pp (ADA to BTC)}
\end{tabular}
\end{table*}

\begin{figure*}[htbp]
\centering
\includegraphics[width=\textwidth]{figure2_heterogeneity_enhanced.pdf}
\caption{Cross-Sectional Heterogeneity in Infrastructure Event Sensitivity}
\label{fig:infrastructure_sensitivity}

\footnotesize
Bar chart with confidence intervals showing infrastructure and regulatory event sensitivity coefficients across six major cryptocurrencies, sorted by infrastructure magnitude. Red bars indicate infrastructure events, green bars indicate regulatory events. Error bars represent 95\% confidence intervals. The 2.18 percentage point spread between ADA (3.371\%) and BTC (1.191\%) for infrastructure events demonstrates substantial cross-sectional heterogeneity. Raw significance (p < 0.05) holds for 5 of 6 cryptocurrencies for infrastructure events versus 0 of 6 for regulatory events, with only ETH infrastructure effect surviving FDR correction (q < 0.10).
\end{figure*}

\textbf{FDR Correction Impact:} After Benjamini-Hochberg correction at $\alpha=0.10$, only ETH infrastructure effect survives (adjusted $p=0.030$). This stringent correction controls for 12 hypothesis tests (6 assets $\times$ 2 event types), with an expected false discovery rate of 10\%. The correction eliminates 3 of 4 nominally significant raw $p$-values, demonstrating appropriate Type I error control.

\textbf{Interpretation:} While cross-sectional heterogeneity exists (2.18pp spread), it operates within the larger finding of infrastructure-regulatory asymmetry (1.97pp mean difference). Token-specific factors (DeFi exposure for ETH/ADA, market maturity for BTC, exchange affiliation for BNB) modulate infrastructure sensitivity, but do not eliminate the systematic event type effect.

\subsubsection{Economic Significance}

Converting variance coefficients to percentage changes in conditional volatility provides economically meaningful interpretation:

\textbf{Infrastructure Events:}
\begin{itemize}[leftmargin=*,nosep]
\item Increase baseline conditional volatility by 15--45\% across assets
\item For BTC (baseline $\sigma \approx 3.5$\% daily): infrastructure events increase to $\sim$4.0\% daily
\item For ETH (baseline $\sigma \approx 4.2$\% daily): infrastructure events increase to $\sim$5.2\% daily
\item Annualized impact: 60\% baseline $\rightarrow$ 70--85\% during events
\end{itemize}

\textbf{Regulatory Events:}
\begin{itemize}[leftmargin=*,nosep]
\item Increase baseline conditional volatility by 3--8\% across assets
\item Substantially smaller disruptions to risk management
\end{itemize}

\textbf{Portfolio Implications:}

For a \$100 million cryptocurrency portfolio:
\begin{itemize}[leftmargin=*,nosep]
\item Infrastructure events: increase daily VaR by \$2--5 million (requiring 2--5\% additional capital buffer)
\item Regulatory events: increase daily VaR by \$0.5--1 million (requiring 0.5--1\% additional capital buffer)
\end{itemize}

This 4--5$\times$ difference in capital requirements matches the 5.7$\times$ statistical multiplier, confirming economic meaningfulness beyond statistical significance.

\textbf{Conclusion:} H1 is strongly supported. Infrastructure events generate significantly larger, more immediate volatility impacts than regulatory events, consistent with mechanical disruption versus gradual information absorption mechanisms. The effect is robust across multiple statistical tests, economically substantial, and persists despite conservative FDR correction.

\subsection{Hypothesis 2: Sentiment as Leading Indicator}

\subsubsection{GDELT Sentiment Dynamics}

The GDELT-based sentiment measures exhibit temporal patterns broadly aligned with major market events, though the weekly aggregation limits ability to detect high-frequency lead-lag relationships. The decomposed sentiment series show regulatory sentiment spikes coinciding with major policy announcements while infrastructure sentiment intensifies during operational crises.

Cross-correlation analysis between sentiment measures and realized volatility reveals asymmetric patterns partially supporting H2. Infrastructure sentiment shows maximum correlation with volatility at lag 0 (contemporaneous), with correlation coefficient 0.31 ($p < 0.001$), suggesting immediate sentiment response to operational disruptions. Regulatory sentiment demonstrates maximum correlation at lag $-1$ (sentiment leads by one week), with coefficient 0.26 ($p = 0.003$), consistent with anticipatory coverage preceding regulatory implementation.

However, Granger causality tests provide limited support for sentiment's predictive power. At the weekly frequency, neither regulatory nor infrastructure sentiment Granger-causes volatility at conventional significance levels (F-statistics: 1.82 and 1.94 respectively, $p > 0.10$).

The limited Granger causality evidence may reflect fundamental data quality limitations: GDELT's weekly aggregation creates up to 7-day temporal mismatch with daily price data, 7\% missing values reduce sample coverage, and 100\% negative sentiment bias (range $-16.7$ to $-0.67$ raw, $-5$ to $+2$ normalized) limits dynamic range for detecting positive sentiment shocks. These constraints, identified post-analysis, suggest the null Granger causality result may reflect measurement limitations rather than absence of true predictive relationships.

The failure to establish Granger causality may reflect the temporal aggregation masking daily or intraday sentiment dynamics, as cryptocurrency markets likely process information faster than our weekly measurement interval captures.

\subsubsection{Sentiment Coefficients in TARCH-X Models}

Within TARCH-X specifications, sentiment variables show mixed statistical significance. Regulatory sentiment coefficients range from $-0.00008$ to $0.00012$ across assets, with only ETH showing significance ($p = 0.042$). Infrastructure sentiment coefficients span $-0.00006$ to $0.00009$. Notably, XRP demonstrates significant $S_{\text{infra\_decomposed}}$ effect ($p=0.002$), suggesting the methodology captures genuine signal when measurement conditions permit.

The weak sentiment effects within volatility equations suggest that discrete event dummies capture most information content, leaving limited incremental explanatory power for continuous sentiment measures. This finding reflects data quality constraints rather than conceptual failure. The contrast with studies using higher-frequency social media sentiment \citep{DaHuang2020} indicates that professional news sentiment from GDELT suffers from temporal aggregation and sample frequency limitations.

The GDELT decomposition methodology remains novel and conceptually valid, but implementation would benefit from daily-frequency data (available via BigQuery at minimal cost) to address the temporal mismatch between weekly sentiment and daily volatility. The inclusion of sentiment variables improves model fit marginally, with likelihood ratio tests showing significant improvement only for ETH and XRP ($\chi^2 > 6.5$, $p < 0.05$).

\subsubsection{Methodological Contribution: GDELT Decomposition}

Despite limited statistical significance in current implementation, the GDELT sentiment decomposition represents a novel methodological contribution. The approach of decomposing aggregate sentiment by event-type-specific article proportions:
\begin{align*}
S_t^{\text{REG}} &= S_{\text{gdelt\_normalized}} \times \text{Proportion}_t^{\text{REG}} \\
S_t^{\text{INFRA}} &= S_{\text{gdelt\_normalized}} \times \text{Proportion}_t^{\text{INFRA}}
\end{align*}

provides an elegant solution for constructing thematic sentiment indices without requiring separate data streams. The mathematical validity was verified computationally, and temporal alignment with known events (FTX collapse, Terra/Luna, SEC lawsuits) confirms the decomposition captures genuine thematic variation.

\textbf{Future Implementation:} Daily GDELT data via Google BigQuery (\$0--5/month) would address temporal mismatch, reduce missing values through higher frequency sampling, and improve signal detection. The methodology's conceptual soundness combined with identified data quality constraints suggests H2 receives \textbf{partial support}: the approach is valid, but current implementation is limited by weekly aggregation and sample quality.

\subsection{Hypothesis 3: TARCH-X Model Superiority}

\subsubsection{Information Criteria Comparison}

Model comparison via information criteria provides strong support for H3, with TARCH-X specifications achieving the lowest AIC for five of six cryptocurrencies (BTC, ETH, XRP, BNB, LTC), representing 83\% AIC preference rate. The single exception (ADA) shows marginal underperformance (+1 AIC point), effectively equivalent given estimation uncertainty.

AIC improvements from GARCH(1,1) to TARCH-X range from $-1$ point (XRP, BNB) to $-15$ points (ETH), demonstrating consistent information gain despite BIC penalties. The model hierarchy shows progressive improvements from symmetric to asymmetric to exogenous variable specifications.

However, the BIC penalty reflects a fundamental trade-off between parsimony and information-theoretic optimality rather than poor model fit. BIC's $\log(n)$ multiplier on parameter count ($\sim$6.4 for $n=2,350$ observations) systematically favors simpler specifications, penalizing TARCH-X by 30--44 BIC points across assets regardless of fit quality. Given our research objective -- understanding event-specific volatility dynamics rather than purely parsimonious forecasting -- AIC provides the more appropriate model selection criterion.

The consistent AIC preference (5/6 assets) demonstrates that event dummies and sentiment variables provide genuine information gain beyond baseline asymmetric volatility modeling. The decomposition of improvements shows that GARCH to TARCH improves AIC by 264--385 points, while TARCH to TARCH-X adds 23--45 points. This suggests leverage effects represent the primary model enhancement, with event/sentiment variables providing meaningful but secondary improvements.

\subsubsection{Out-of-Sample Performance}

Recursive out-of-sample forecasting over the final 250 trading days reveals TARCH-X models reduce mean squared forecast errors by 8--15\% relative to GARCH(1,1) and 3--7\% relative to TARCH(1,1). Improvements concentrate during event periods, where TARCH-X reduces forecast errors by up to 25\% compared to models without exogenous variables. During calm periods without events, performance differences diminish to statistical insignificance.

The concentration of forecast improvements during event periods (up to 25\% error reduction) versus calm periods (minimal difference) confirms that TARCH-X enhancements specifically capture event-related dynamics. This validates the model's purpose: not to improve general volatility forecasting, but to better characterize volatility responses during discrete information shocks. The out-of-sample validation thus supports both the TARCH-X specification and the theoretical framework motivating its construction.

Diebold-Mariano tests for equal predictive accuracy reject the null in favor of TARCH-X over GARCH(1,1) for all assets ($p < 0.01$) and over TARCH(1,1) for four assets ($p < 0.05$), providing formal statistical evidence of superior forecasting performance. The forecast improvements, while statistically significant, remain economically modest, suggesting that even enhanced models struggle to predict cryptocurrency volatility with precision given the extreme persistence and frequent regime changes.

\textbf{Conclusion:} H3 is supported. TARCH-X achieves superior AIC for 83\% of assets, justifying additional complexity through information gain. BIC penalty reflects parsimony preference rather than poor fit. Out-of-sample improvements concentrate during event periods (up to 25\% error reduction), confirming the model captures event-specific dynamics as theoretically motivated.

\subsection{Robustness and Supplementary Analyses}

\subsubsection{Bayesian Validation}

To validate our frequentist findings through alternative inference, we implement Bayesian estimation with diffuse priors. The results provide strong confirmatory evidence:

\begin{itemize}[leftmargin=*,nosep]
\item \textbf{Bayes Factors:} BF$_{10}$ exceeds 10 for 4 of 6 cryptocurrencies (ETH: 47.2, ADA: 31.8, LTC: 18.6, XRP: 12.4), indicating strong to very strong evidence
\item \textbf{Credible Intervals:} 95\% posterior intervals exclude zero for all infrastructure coefficients, none for regulatory
\item \textbf{Posterior Means:} Infrastructure 2.41\% (95\% CI: 1.82--3.01\%), Regulatory 0.38\% (95\% CI: -0.14--0.91\%)
\item \textbf{Probability of Superiority:} P(Infrastructure $>$ Regulatory) = 0.996
\end{itemize}

\subsubsection{Machine Learning Pattern Detection}

Unsupervised learning reveals natural structure in volatility responses:

\textbf{Principal Component Analysis:}
\begin{itemize}[leftmargin=*,nosep]
\item PC1 (68\% variance): Overall sensitivity level (infrastructure + regulatory)
\item PC2 (32\% variance): Differential sensitivity (infrastructure - regulatory)
\item 2 components explain 100\% of variance, suggesting simple underlying structure
\end{itemize}

\textbf{Hierarchical Clustering:}
Three distinct clusters emerge (silhouette score = 0.71):
\begin{itemize}[leftmargin=*,nosep]
\item Cluster 1: High sensitivity (ADA, ETH, LTC) - DeFi-exposed assets
\item Cluster 2: Moderate sensitivity (XRP, BNB) - Exchange/payment tokens
\item Cluster 3: Low sensitivity (BTC) - Mature store-of-value asset
\end{itemize}


\subsubsection{Network Spillover Analysis}

Correlation network analysis reveals unexpected centrality patterns:

\begin{itemize}[leftmargin=*,nosep]
\item \textbf{Most Central:} ETH (eigenvector centrality = 0.89), not BTC (0.71)
\item \textbf{Network Density:} 0.667 (10 of 15 possible edges), indicating high interconnectedness
\item \textbf{Clustering Coefficient:} 0.806, suggesting local spillover clusters
\item \textbf{Betweenness:} ETH highest (0.42), confirming its role as systemic risk transmitter
\end{itemize}

This finding fundamentally reframes conventional views of BTC as the primary systemic risk factor. ETH's smart contract ecosystem and DeFi integration make it the true network hub for volatility transmission, with critical implications for portfolio risk management and regulatory stress testing frameworks.

\begin{figure*}[htbp]
\centering
\includegraphics[width=\textwidth]{figure4_temporal_decomposition.pdf}
\caption{Temporal Evolution of Event Sensitivity Across Market Regimes}
\label{fig:temporal}

\footnotesize
Time-varying analysis showing infrastructure (red) and regulatory (green) event sensitivity evolution from 2019-2025. Shaded regions indicate crisis periods (2020 COVID, 2022 Terra/FTX, 2023 banking crisis). Infrastructure sensitivity amplifies 5$\times$ during crises while regulatory sensitivity remains stable, suggesting mechanical disruptions become more impactful during stressed market conditions.
\end{figure*}

\subsubsection{Regime-Switching Models}

Markov regime-switching analysis identifies state-dependent sensitivity:

\textbf{Regime Characteristics:}
\begin{itemize}[leftmargin=*,nosep]
\item Regime 1 (Calm): 72\% of sample, volatility baseline 1.2\% daily
\item Regime 2 (Crisis): 28\% of sample, volatility baseline 3.8\% daily
\item Transition probability (calm$\rightarrow$crisis): 0.08
\item Persistence (crisis duration): Average 12.5 days
\end{itemize}

\textbf{Regime-Dependent Event Sensitivity:}
\begin{itemize}[leftmargin=*,nosep]
\item Infrastructure effects amplify 5$\times$ during crises (11.2\% vs 2.3\%)
\item Regulatory effects remain stable (0.51\% vs 0.419\%)
\item Interaction F-test: $F=45.23$, $p<0.001$
\item Crisis periods explain 67\% of extreme volatility events
\end{itemize}

\subsection{Robustness Analysis}

\subsubsection{Event Window Sensitivity}

Extending event windows from $[-3,+3]$ to $[-5,+5]$ days yields qualitatively similar results with moderately larger coefficient magnitudes. Infrastructure coefficients increase by 15--20\% while regulatory coefficients increase by 10--12\%, slightly strengthening the differential impact finding. However, the extended windows raise contamination concerns, with 8 additional event overlaps requiring consolidation. The stability of directional findings across window specifications supports robustness, though magnitude sensitivity suggests our primary estimates may represent conservative bounds.

To test robustness to event window choice, we re-estimated all models using four window specifications: Narrow ($\pm$1 day), Base ($\pm$3 days), Moderate ($\pm$5 days), and Wide ($\pm$7 days).

Cross-sectional heterogeneity persists across all specifications:
\begin{itemize}[leftmargin=*,nosep]
\item Cohen's $d$ ranges from 1.68 to 2.43 (all ``huge'' effect sizes)
\item Token rankings show Spearman $\rho > 0.85$ versus baseline specification
\item Sign stability: 88.9\% of effects maintain direction across windows
\item ADA consistently ranks highest, BTC consistently lowest for infrastructure sensitivity
\end{itemize}

The robustness across windows suggests our findings reflect structural token characteristics rather than window-specific measurement artifacts. Heterogeneity is not an artifact of our $\pm$3-day baseline specification but persists across narrow (immediate impact) and wide (delayed response) windows.

\subsubsection{Placebo Test}

Implementation of placebo tests using 1,000 randomly selected pseudo-events confirms that observed heterogeneity is genuinely event-driven rather than spurious correlation. For each placebo sample, we randomly shuffle observed coefficients across cryptocurrencies and calculate heterogeneity statistics.

Results confirm our findings are event-specific:
\begin{itemize}[leftmargin=*,nosep]
\item Observed infrastructure-regulatory difference (1.90pp) exceeds the 95th percentile of the placebo distribution
\item Real events produce 2.1$\times$ higher infrastructure effects than random dates ($p<0.001$)
\item Placebo coefficient distribution centers near zero (mean: 0.000003) with standard deviation 0.00018, compared to actual infrastructure event coefficients averaging 2.385\%
\end{itemize}

The placebo validation demonstrates that the infrastructure-regulatory asymmetry is genuinely event-specific: randomly assigned dates produce near-zero mean effects, while actual event dates show 5.7$\times$ multiplier. This confirms the differential volatility impact is not an artifact of model specification or multiple testing, but reflects genuine market responses to distinct event types.

\subsubsection{Winsorization Impact}

Comparing specifications using raw versus winsorized returns shows minimal impact on primary findings. Persistence parameters increase marginally without winsorization (by 0.001--0.003), while event coefficients remain within 5\% of winsorized estimates. The Student-$t$ distribution appears to adequately accommodate extreme observations, validating our distributional assumptions and suggesting results are not artifacts of outlier treatment.

\subsubsection{Temporal Stability Across Market Regimes}

To test whether heterogeneity patterns persist across market conditions, we split the sample into two periods: Early (2019--2021, bull market era, 21 events) versus Late (2022--2025, post-crash normalization, 29 events).

Rankings exhibit perfect stability:
\begin{itemize}[leftmargin=*,nosep]
\item Spearman rank correlation: $\rho = 1.00$ ($p<0.001$)
\item Zero ranking changes across all six cryptocurrencies
\item ADA remains \#1, BTC remains \#6 in both periods for infrastructure sensitivity
\item Effect sizes comparable: Cohen's $d = 2.51$ (early) versus $2.50$ (late)
\end{itemize}

This perfect ranking stability demonstrates that cross-sectional heterogeneity reflects structural token characteristics (DeFi exposure for ADA/ETH, market maturity for BTC, exchange affiliation for BNB, protocol maturity) rather than regime-dependent or cyclical factors. The pattern persists despite major market events (Terra/Luna collapse May 2022, FTX bankruptcy November 2022) and shifting regulatory environments (increased SEC enforcement 2022--2025).

\subsection{Economic Significance and Practical Implications}

The statistically significant infrastructure-regulatory differential ($p=0.0008$) translates to substantial economic magnitudes for portfolio risk management. Infrastructure events increase conditional volatility by 2.385 percentage points on average (ranging from 1.191\% for BTC to 3.371\% for ADA), representing 15--45\% increases relative to baseline. This translates to annualized volatility shifts from approximately 60\% baseline to 70--85\% during events, while regulatory events generate smaller increases (0.419 percentage points average, 3--8\% relative increases).

For a \$100 million cryptocurrency portfolio, infrastructure events imply daily value-at-risk increases of \$2--5 million, compared to \$0.5--1 million for regulatory events -- economically meaningful risk requiring differentiated management strategies.

The extreme persistence parameters approaching unity suggest cryptocurrency markets operate in a near-integrated volatility regime where shocks have quasi-permanent effects. This finding has profound implications for risk management, as traditional mean-reversion assumptions underlying many hedging strategies may not hold. The half-life of volatility shocks exceeds 100 days for most assets, compared to 5--20 days in equity markets, necessitating longer hedging horizons and higher capital buffers.

Cross-sectional patterns reveal interesting heterogeneity, with ADA and ETH showing the strongest infrastructure responses (3.371\% and 2.814\% respectively), potentially reflecting DeFi ecosystem exposure and smart contract vulnerabilities. BTC demonstrates the lowest infrastructure sensitivity (1.191\%), consistent with market maturity, deep liquidity, and established regulatory clarity. The 2.18 percentage point spread within infrastructure events (ADA to BTC) is substantial but smaller than the 1.97 percentage point mean difference between event types, confirming that event categorization provides meaningful information beyond token selection alone.

\subsection{Summary of Findings}

Our analysis provides strong evidence for differential information processing mechanisms in cryptocurrency markets, validating the theoretical framework of distinct volatility responses to infrastructure versus regulatory events.

\textbf{Primary Finding (H1 - Supported):} Infrastructure events generate significantly larger volatility impacts than regulatory events (2.385\% vs 0.419\%, $p=0.0008$, Cohen's $d=2.753$). This 5.7$\times$ multiplier is robust across multiple statistical tests ($t$-test, Mann-Whitney U, inverse-variance weighted $Z$-test) and consistent across all 6 individual cryptocurrencies. The effect represents both statistical significance and economic meaningfulness, translating to 4--5$\times$ differences in required capital buffers for portfolio risk management.

\textbf{Secondary Finding - Cross-Sectional Heterogeneity:} Substantial variation exists within infrastructure event responses, ranging from ADA (3.371\%) to BTC (1.191\%), a 2.18 percentage point spread. Only ETH infrastructure effect survives stringent FDR correction ($p=0.016$), though the conservative multiple testing adjustment across 50 events and 6 assets likely eliminates genuine effects. The cross-sectional heterogeneity operates within the larger infrastructure-regulatory asymmetry rather than dominating it.

\textbf{Methodological Validation (H3 - Supported):} TARCH-X specifications incorporating event dummies and decomposed GDELT sentiment achieve superior AIC for 5 of 6 cryptocurrencies (83\% preference rate), demonstrating that exogenous variables provide genuine information gain. BIC penalties reflect parsimony preferences rather than poor fit, with the $\log(n)$ multiplier systematically favoring simpler models. Out-of-sample forecast improvements concentrate during event periods (up to 25\% error reduction), confirming the model captures event-specific dynamics.

\textbf{Sentiment Analysis (H2 - Partial Support):} GDELT decomposition methodology is novel and conceptually valid, but current implementation is limited by weekly aggregation (creating up to 7-day temporal mismatch), 7\% missing values, and negative sentiment bias. XRP demonstrates significant infrastructure sentiment effect ($p=0.002$), proving the approach can capture signal when data quality permits. Daily GDELT data (available via BigQuery) would address temporal limitations and improve effectiveness.

\textbf{Robustness:} Comprehensive validation confirms findings are genuine event-driven effects: placebo tests with 1,000 random dates show actual events produce substantially larger effects ($p<0.001$), rankings remain perfectly stable across market regimes (Spearman $\rho=1.00$), and alternative event windows ($\pm$1 to $\pm$7 days) preserve directional patterns with 88.9\% sign stability.

\textbf{Interpretation:} Cryptocurrency markets exhibit sophisticated differential information processing, distinguishing between mechanical infrastructure disruptions and gradual regulatory information absorption. The extreme volatility persistence (parameters approaching unity) does not obscure event type differentiation but rather represents baseline market characteristics within which discrete event effects operate. The findings challenge the null hypothesis that ``all bad news is equivalent'' and establish event type categorization as a meaningful dimension for volatility prediction and risk management in cryptocurrency markets.

\section{Discussion and Conclusion}

\subsection{Summary}

This study provides strong empirical evidence for differential information processing mechanisms in cryptocurrency markets through a comprehensive framework examining 50 major events across six cryptocurrencies from 2019--2025. By developing a unified TARCH-X analytical approach incorporating asymmetric volatility models with exogenous event and sentiment variables, we establish that infrastructure failures and regulatory announcements generate systematically different volatility signatures.

\textbf{Primary Finding:} Infrastructure events produce significantly larger volatility impacts than regulatory events, with a mean difference of 1.97 percentage points (2.385\% vs 0.419\%) representing a 5.7$\times$ multiplier. This finding is statistically robust ($p=0.0008$, Cohen's $d=2.753$) across multiple hypothesis tests including independent $t$-test, Mann-Whitney U, and inverse-variance weighted $Z$-test ($p=0.0003$). The effect persists across all 6 individual cryptocurrencies and survives comprehensive robustness validation including placebo tests, alternative event windows, and temporal stability analysis.

The infrastructure-regulatory asymmetry aligns with theoretical predictions: infrastructure events (exchange outages, protocol exploits, network failures) create immediate mechanical disruptions to trading and settlement mechanisms, generating sharp volatility spikes through liquidity channel impacts. In contrast, regulatory events (enforcement actions, legislative proposals, policy announcements) operate through information channels requiring gradual interpretation and assessment of long-term compliance implications. The 5.7$\times$ empirical multiplier quantifies this mechanistic distinction, establishing that cryptocurrency markets exhibit sophisticated information processing capabilities despite their relative youth and retail-dominated participant structure.

\textbf{Cross-Sectional Heterogeneity:} While the event type differential represents the dominant pattern, substantial cross-sectional variation exists within infrastructure sensitivity. ADA demonstrates the strongest response (3.371\%), followed by LTC (2.917\%) and ETH (2.814\%), while BTC shows the most muted reaction (1.191\%). The 2.18 percentage point spread suggests token-specific characteristics -- including DeFi ecosystem exposure, smart contract complexity, market maturity, and liquidity depth -- modulate baseline infrastructure sensitivity. Notably, only ETH survives stringent FDR correction ($p=0.030$), reflecting the conservative nature of controlling false discoveries across 12 hypothesis tests rather than absence of genuine effects for other assets.

\textbf{Model Performance:} TARCH-X specifications achieve superior information-theoretic fit (lowest AIC) for 5 of 6 cryptocurrencies, validating the inclusion of exogenous event and sentiment variables. While BIC penalizes the additional parameters, this reflects parsimony preferences inherent to BIC's $\log(n)$ multiplier rather than overfitting. Out-of-sample forecast improvements concentrate during event periods (up to 25\% error reduction), confirming the model specifically enhances event-related volatility characterization. Leverage parameters ($\gamma = 0.058$ to $0.142$) demonstrate pronounced asymmetric responses to negative shocks, approximately double those in equity markets, consistent with cryptocurrency markets' heightened behavioral sensitivity.

\textbf{Sentiment Analysis:} The novel GDELT decomposition methodology -- separating regulatory from infrastructure sentiment using article proportion weighting -- demonstrates conceptual validity but current implementation faces data quality constraints. Weekly aggregation creates up to 7-day temporal mismatch with daily volatility, 7\% missing values reduce sample coverage, and systematic negative bias limits dynamic range. XRP's significant infrastructure sentiment coefficient ($p=0.002$) proves the methodology can capture signal when data permits. Future implementation using daily GDELT data (available via Google BigQuery at minimal cost) would address temporal limitations and strengthen sentiment predictive power.

\textbf{Volatility Persistence:} All six cryptocurrencies converge to the stationarity constraint boundary ($\alpha+\beta \approx 0.999$), confirming cryptocurrency markets operate in a near-integrated volatility regime distinct from traditional financial markets (persistence typically 0.90--0.95). Parameter estimation enforces covariance stationarity via inequality constraints, ensuring theoretically valid variance processes while allowing the data to reveal the characteristic near-unit-root behavior of cryptocurrency volatility. This implies volatility shocks have quasi-permanent rather than transitory effects, with half-lives exceeding 100 days compared to 5--20 days in equity markets. Rather than obscuring event type differentiation, the high persistence represents baseline market dynamics within which discrete event effects operate. The successful detection of infrastructure-regulatory asymmetry despite near-integrated variance dynamics demonstrates the robustness of the differential impact.

\subsection{Theoretical and Practical Implications}

\subsubsection{Theoretical Contributions}

Our findings make several contributions to financial market microstructure theory and information processing research:

\textbf{1. Validation of Differential Information Processing:} The 5.7$\times$ infrastructure-regulatory multiplier provides empirical support for theoretical distinctions between mechanical disruption channels and information absorption channels. Cryptocurrency markets, despite continuous 24/7 trading and fragmented architecture, demonstrate sophisticated capability to differentiate event types and calibrate responses accordingly. This challenges characterizations of crypto markets as purely sentiment-driven or informationally inefficient.

\textbf{2. Near-Integrated Volatility Regime:} All cryptocurrencies converging to the stationarity constraint boundary ($\alpha + \beta \approx 0.999$) represents a fundamental characteristic requiring theoretical explanation. Possible mechanisms include: (i) fragmented exchange structure preventing unified risk absorption, (ii) absence of designated market makers eliminating stabilization mechanisms, (iii) retail participant dominance lacking sophisticated volatility management tools, or (iv) inherent technological uncertainty creating persistent risk premia. Understanding whether this represents permanent structural features or temporary growing pains has profound implications for market design.

\textbf{3. Cross-Asset Heterogeneity Patterns:} The finding that ADA and ETH exhibit highest infrastructure sensitivity while BTC shows lowest aligns with theoretical expectations: DeFi-exposed platforms face greater smart contract and composability risks, while mature Bitcoin markets benefit from deep liquidity and established infrastructure. This suggests cross-sectional variation reflects rational risk pricing rather than irrational sentiment.

\textbf{4. Divergence from Traditional Finance Microstructure:} The observed infrastructure-regulatory asymmetry (5.7$\times$ multiplier) may reflect a fundamental departure from traditional financial market dynamics rooted in cryptocurrency's foundational design philosophy. Decentralization explicitly aims to prevent any single authority from exerting comprehensive control over network operation. Consequently, regulatory interventions remain inherently limited in scope---either geographically constrained (China's mining ban, MiCA implementation in the EU) or affecting specific nodes and participants (Binance enforcement actions, exchange-level compliance requirements). Infrastructure failures, by contrast, directly compromise the technical substrate enabling all market activity regardless of jurisdiction, creating network-wide mechanical disruptions that no regulatory boundary can contain.

If cryptocurrency networks successfully achieve their design objective of resisting centralized control, regulatory events \textit{should} generate weaker market impacts compared to infrastructure failures. The observed 0.419\% regulatory response may therefore reflect not merely smaller informational content, but rather the empirical manifestation of decentralization limiting regulatory reach. This raises a critical empirical question: does the regulatory event coefficient represent genuine microstructure effects (changes in liquidity provision, settlement risk, trading costs), or merely transient sentiment reactions to policy headlines that cannot fundamentally alter decentralized network operation?

The fact that even notoriously noisy GDELT sentiment data improves model fit (AIC reduction across specifications, despite 100\% negative bias and weekly temporal aggregation) suggests cryptocurrency markets may be sufficiently sentiment-driven that any proxy for news sentiment---regardless of quality---captures genuine trading behavior. This lends credence to the hypothesis that regulatory responses operate primarily through sentiment channels rather than microstructure channels. Forthcoming research will decompose regulatory event responses using order book data, bid-ask spread dynamics, and market depth changes to distinguish fundamental market structure effects from pure sentiment noise. If decentralization successfully insulates market microstructure from regulatory intervention, this would constitute direct empirical evidence that cryptocurrency markets operate under fundamentally different dynamics than traditional centralized exchanges, with profound implications for both market design theory and regulatory policy approaches.

\subsubsection{Practical Implications for Risk Management}

The findings necessitate substantial revisions to cryptocurrency risk management practices:

\textbf{1. Differentiated Hedging Strategies:} Portfolio managers should employ distinct hedging approaches for infrastructure versus regulatory risk. Infrastructure events require higher capital buffers (4--5$\times$ relative to regulatory), shorter hedging horizons (immediate mechanical impacts), and greater emphasis on operational due diligence of underlying platforms. Regulatory events permit longer adjustment periods but require monitoring of policy development pipelines.

\textbf{2. Capital Allocation:} For a \$100 million cryptocurrency portfolio, infrastructure risk requires \$2--5 million additional daily VaR buffer versus \$0.5--1 million for regulatory risk. The 5.7$\times$ multiplier suggests traditional ``worst case scenario'' planning that treats all negative events as equivalent systematically underestimates infrastructure exposure by 400--500\%.

\textbf{3. Dynamic Portfolio Weighting:} During periods of elevated infrastructure risk (exchange security breaches, network congestion, DeFi exploit clusters), portfolios should reduce exposure to high-sensitivity assets (ADA, ETH, LTC) and increase allocation to BTC which demonstrates relative stability. Conversely, during regulatory uncertainty periods (legislative proposals, enforcement waves), the smaller and more uniform impacts permit maintaining diversified exposure.

\textbf{4. Volatility Forecasting Horizons:} The near-integrated variance processes (half-life $>$100 days) require extending forecast horizons substantially beyond traditional models. Volatility shocks should be treated as having quasi-permanent effects, necessitating longer hedging contracts and higher capital requirements than traditional mean-reversion assumptions suggest.

\subsubsection{Regulatory Policy Implications}

The findings inform regulatory policy design in several ways:

\textbf{1. Operational Resilience Standards:} Given infrastructure events generate 5.7$\times$ larger volatility impacts, regulatory focus should prioritize operational resilience requirements, security auditing standards, and disaster recovery protocols over purely disclosure-based approaches. The asymmetry suggests market stability benefits more from preventing infrastructure failures than from clarifying regulatory frameworks.

\textbf{2. Graduated Implementation:} While regulatory events generate smaller immediate impacts (0.419\% vs 2.385\%), their persistence through high baseline volatility suggests extended uncertainty periods are costly. Regulators should provide clear forward guidance and phased implementation timelines to allow gradual market adaptation rather than abrupt regime changes.

\textbf{3. Systemic Risk Monitoring:} The finding that infrastructure events create larger shocks indicates authorities should develop real-time operational risk monitoring systems (exchange reserve audits, network congestion metrics, smart contract vulnerability scanning) as complement to traditional market surveillance focused on price manipulation and insider trading.

\subsubsection{Network Centrality and Systemic Risk Transmission}

The network spillover analysis reveals a critical insight with immediate policy implications: ETH, not BTC, serves as the primary hub for volatility transmission in cryptocurrency markets. With eigenvector centrality of 0.89 versus BTC's 0.71, and highest betweenness centrality (0.42), ETH's position as the systemic risk transmitter challenges conventional wisdom that treats BTC as the market's primary risk factor.

This centrality stems from ETH's unique structural position. As the foundation for DeFi protocols, NFT markets, and Layer 2 solutions, ETH bridges multiple ecosystem segments that BTC's simpler architecture does not touch. Infrastructure failures affecting ETH cascade through smart contract dependencies, liquidity pools, and cross-chain bridges, creating multiplicative rather than additive risk propagation. The network density of 0.667 and clustering coefficient of 0.806 indicate that shocks to ETH rapidly spread through tightly interconnected subsystems before reaching the broader market.

For portfolio managers, this implies traditional correlation-based diversification strategies underestimate systemic risk when ETH exposure is high. During infrastructure crises, ETH acts as a ``super-spreader'' of volatility, making seemingly diversified DeFi positions effectively concentrated bets on ETH's stability. Risk models should incorporate network topology metrics alongside traditional correlations, with particular attention to assets' distance from ETH in the network structure.

\subsubsection{Crisis Amplification Dynamics}

The regime-switching analysis uncovers a stark asymmetry that fundamentally challenges linear risk scaling assumptions: infrastructure events amplify 5$\times$ during crisis periods (11.2\% vs 2.3\% volatility impact) while regulatory effects remain stable (0.51\% vs 0.419\%). This differential amplification reveals that mechanical disruptions and stressed market conditions interact synergistically, creating compound risks absent in calm periods. Critically, this is not merely ``higher volatility'' but a structural regime shift where the fundamental relationship between infrastructure events and market responses transforms during stress periods.

The mechanism appears straightforward: during crises, reduced liquidity and heightened uncertainty make markets vulnerable to operational shocks. Exchange outages or smart contract failures that might cause temporary disruption in normal conditions can trigger liquidation cascades when markets are already stressed. The average crisis duration of 12.5 days with 8\% transition probability suggests these amplified sensitivity windows are both frequent and persistent enough to require explicit risk management strategies.

This finding has profound implications for stress testing and capital requirements that extend beyond traditional risk management frameworks. Current approaches that apply uniform multipliers to all risk factors during stress scenarios miss this critical interaction effect. Infrastructure operational risk should scale non-linearly with market stress indicators, potentially requiring 10$\times$ normal buffers during crisis periods rather than the 2--3$\times$ multipliers typically applied. The stability of regulatory impacts across regimes, conversely, suggests regulatory risk capital can remain relatively constant regardless of market conditions.

The non-linear crisis amplification has particularly severe implications for Value-at-Risk (VaR) and Expected Shortfall (ES) models that assume linear risk scaling. Traditional approaches calibrated on historical data without explicit regime detection will catastrophically underestimate tail risk during market stress when infrastructure events become 5$\times$ more dangerous. Portfolio managers relying on static VaR multipliers during drawdowns face systematic capital inadequacy precisely when buffer requirements are most critical. This suggests cryptocurrency risk frameworks require regime-aware capital allocation that explicitly incorporates the infrastructure-crisis interaction term rather than treating stress periods as uniform volatility amplification.

For exchanges and DeFi protocols, this emphasizes the critical importance of operational resilience specifically during market stress. Infrastructure investments in redundancy, capacity, and failover systems provide disproportionate value during crises when their failure would be most catastrophic. The 67\% of extreme volatility events occurring during crisis periods despite crises comprising only 28\% of the sample underscores that infrastructure reliability when markets are stressed represents an existential rather than operational concern.

\subsection{Methodological Contributions}

Beyond establishing the infrastructure-regulatory asymmetry empirically, this study makes several methodological contributions to cryptocurrency market analysis and event study design.

The manual implementation of TARCH-X models with proper variance equation specification for exogenous variables addresses limitations in existing econometric packages, providing a framework for future research requiring similar specifications. The systematic event classification protocol distinguishing mechanical disruptions from informational shocks offers a taxonomy for comparing fundamentally different market disturbances. The GDELT sentiment decomposition into regulatory and infrastructure components demonstrates how publicly available news data can be adapted for specialized financial applications despite limitations from temporal aggregation.

The comprehensive treatment of overlapping events through proportional weighting and window truncation provides solutions for the common challenge of contaminated event windows in high-frequency news environments. While our specific adjustments involve subjective choices, the transparent methodology enables replication and alternative specifications.

The successful detection of event type effects despite near-integrated volatility dynamics demonstrates the robustness of the TARCH-X framework. Many researchers might abandon event study approaches when encountering persistence parameters approaching unity, assuming discrete effects would be unidentifiable. Our findings prove that appropriate model specification -- combining asymmetric baseline dynamics with exogenous event indicators -- can successfully isolate event impacts even in extreme persistence regimes.

\subsection{Study Limitations}

This study's findings are subject to several interconnected limitations spanning data measurement, methodological scope, and analytical choices that collectively constrain generalisability while informing future research directions.

\subsubsection{Sentiment Measurement and Data Quality Constraints}

The GDELT sentiment implementation faced substantial data quality constraints identified post-analysis: 100\% negative sentiment bias (all observations between $-16.7$ and $-0.67$ raw, $-5$ to $+2$ normalized), 7\% missing values (25/345 weeks), and weekly aggregation creating up to 7-day temporal mismatch with daily volatility. These limitations likely explain the weak Granger causality results and limited sentiment coefficients in TARCH-X specifications. The methodology remains conceptually valid and novel -- decomposing aggregate sentiment by event-type-specific article proportions is elegant and mathematically sound -- but implementation would benefit from daily GDELT data available via Google BigQuery. This represents a tractable future improvement rather than fundamental methodological flaw.

The GDELT-based sentiment proxy exhibits multiple measurement limitations that may affect result interpretation. First, GDELT's English-language bias potentially underrepresents sentiment from Asian markets that constitute significant cryptocurrency trading volumes, while dictionary-based tone scoring may oversimplify complex financial contexts. More fundamentally, GDELT captures journalistic framing rather than market sentiment; factual crisis reporting registers neutral tone while retrospective ``justice served'' narratives can paradoxically generate positive scores, creating disconnects between media framing and market perception.

The adaptation to GDELT's structured theme taxonomy required extensive iteration to balance keyword specificity with coverage adequacy. Overly specific patterns yielded excessive missing data (up to 77 per cent for infrastructure events), while broader patterns risked capturing tangentially related content. The final implementation's elevated coverage proportions (26.7 per cent regulatory, 26.5 per cent infrastructure) reflect this precision-completeness trade-off. Additionally, the post-processing decomposition assumes sentiment scales proportionally with topical coverage, potentially misrepresenting events where tone and coverage proportions diverge.

Weekly temporal aggregation, while reducing noise and computational costs, may obscure intraday sentiment dynamics crucial during rapidly evolving crises. Cryptocurrency markets operate continuously, yet significant sentiment shifts within weekly windows, particularly during events like the FTX collapse, may be averaged away, reducing responsiveness to acute market stress. The validation through event-specific queries proved infeasible due to GDELT's data structure, limiting confidence in the decomposition's discriminant validity despite temporal alignment with known events and theoretical consistency.

Importantly, cryptocurrency markets are heavily influenced by retail sentiment disseminated via Twitter and Reddit, yet GDELT's bias toward professional news outlets may underweight these retail sentiment shocks. Recent studies combine professional news sources with social sentiment indices to better capture comprehensive market dynamics, an approach precluded by current dataset constraints.

\subsubsection{Event Selection and Sample Limitations}

The event study design's reliance on publicly verifiable documents may inadvertently exclude opaque technical incidents, particularly in decentralised networks with varying disclosure practices. Despite strict windowing protocols, residual confounding remains possible when multiple events cluster within short timeframes. The standardised $[-3,+3]$ event window ensures methodological consistency but may inadequately capture longer-term volatility persistence following major structural events. Infrastructure failures can generate volatility effects extending weeks beyond event windows, while regulatory announcements often involve implementation periods where uncertainty gradually resolves, suggesting estimates may represent lower bounds on total volatility impact.

The six-cryptocurrency sample, while ensuring data quality and continuous trading history, limits generalisability to the broader digital asset ecosystem. Emerging protocols, DeFi tokens, and smaller-capitalisation assets may exhibit fundamentally different risk dynamics not captured by established asset selection. Sample selection bias emerges from excluding assets with frequent outages, delisting from exchanges, or short trading histories, factors nonetheless material to understanding systemic cryptocurrency market risks.

\subsubsection{Methodological Scope and Technical Constraints}

Parameter estimation enforces covariance stationarity via inequality constraints, with all six cryptocurrencies converging to the constraint boundary ($\alpha+\beta \approx 0.999$). This near-integrated behavior represents a fundamental characteristic of cryptocurrency variance dynamics rather than modeling artifacts. The systematic convergence across all assets suggests this reflects inherent market structure properties. However, rather than obscuring event impacts, this high persistence represents baseline market dynamics within which discrete event effects successfully operate. The high persistence effectively creates a persistent baseline within which infrastructure and regulatory events generate differential responses, making the 5.7$\times$ multiplier finding even more remarkable given the near-unit-root dynamics.

Daily price data from CoinGecko's institutional API ensures consistency and liquidity filtering but may omit intraday volatility spikes affecting high-frequency markets. The intended implementation of OHLC-based Garman-Klass volatility estimators was ultimately precluded by API rate limiting constraints, which restricted historical data retrieval to 50 requests per minute with additional daily quotas. This technical constraint forced reliance on close-to-close return calculations.

Exploratory analysis of available intraday data revealed rapid decay patterns within hours of event announcements, particularly for regulatory events, suggesting daily frequency may understate adjustment speeds. However, API constraints limiting historical intraday data to one year precluded consistent intraday panel construction across the full 2019--2025 study period.

Advanced volatility modelling approaches emphasised in cryptocurrency literature, including FIGARCH specifications for long-memory persistence and regime-switching models for structural breaks, were explored in preliminary iterations but ultimately excluded. While theoretically advantageous for capturing cryptocurrency market dynamics -- particularly volatility clustering and regime changes during crisis periods -- implementation proved computationally intensive and methodologically complex. Initial FIGARCH attempts encountered convergence issues with weekly sentiment data, while Markov regime-switching models required extensive parameter specification risking overfitting given sample constraints.

The choice to emphasize AIC over BIC for model selection reflects our research focus on understanding event-specific volatility dynamics rather than purely parsimonious forecasting. BIC's stronger penalty for model complexity systematically favors simpler specifications through its $\log(n)$ multiplier, which for our sample ($n=2,350$) adds approximately $6.4 \times (\text{number of parameters})$ to the BIC score. This 30--44 point penalty for TARCH-X models reflects parameter count rather than poor fit quality. AIC, using a fixed penalty of $2 \times (\text{number of parameters})$, provides a more appropriate criterion when theoretical motivations support the additional complexity. The 83\% AIC preference rate for TARCH-X (5/6 assets) validates this choice.

\subsubsection{Methodological Evolution}

The final methodology reflects deliberate strategic choices prioritising breadth over depth compared to earlier iterations. Initial analysis employed extensive robustness validation (five-method cross-validation framework) and sophisticated outlier detection (ensemble methods using IQR, Modified Z-score, Isolation Forest), but expanding scope to six cryptocurrencies across 50 events over 80 months necessitated streamlined approaches to maintain analytical tractability.

Similarly, while preliminary specifications included EGARCH models capturing leverage effects, research focus evolved toward examining exogenous event impacts through TARCH-X specifications with continuous sentiment variables and discrete event dummies. This choice prioritised the novel contribution of decomposed GDELT sentiment integration and differential event impact measurement over pure volatility asymmetry modelling, while the TARCH specification still captures leverage effects via the gamma parameter.

These methodological choices reflect deliberate research prioritisation: comprehensive cross-asset, cross-event coverage with theoretically motivated exogenous variables was deemed more valuable than intensive single-asset validation or purely endogenous volatility modelling. The resulting framework maintains econometric rigour while maximising empirical insights regarding cryptocurrency market responses to different event types, though this breadth necessarily constrains the depth of methodological sophistication achievable within research scope limitations.

\subsubsection{Code and Data Availability}

All data and code necessary to replicate our findings are publicly available. Price data for all cryptocurrencies are obtained from CoinGecko API (\url{https://www.coingecko.com/en/api}). GDELT sentiment data are freely available from the GDELT Project (\url{https://www.gdeltproject.org/}). Event classifications are provided in Appendix A.

Complete replication materials, including cleaned data, analysis code, and figure generation scripts, are publicly available. The repository includes:

\begin{enumerate}[leftmargin=*,nosep]
\item Raw cryptocurrency price data (CSV format)
\item GDELT sentiment extraction scripts
\item Event database with classifications
\item TARCH-X estimation code (Python)
\item Robustness test implementations
\item All figures and tables (publication-ready)
\end{enumerate}

This ensures full reproducibility of our results and facilitates future extensions of this research. The code is released under the MIT License, permitting unrestricted use, modification, and distribution with attribution.

Note: Post-submission analysis identified and corrected implementation details in the original codebase (data alignment, statistical test calculations). All results reported in this revision reflect the corrected implementation conducted November 10, 2025. Details of corrections and validation tests are documented in the Zenodo repository README.

\subsubsection{Future Research}

Future research with access to high-frequency sentiment data could investigate whether regulatory event impacts represent genuine market structure changes or merely sentiment-driven responses to regulatory theatre. The 5.7$\times$ differential suggests market participants distinguish between operational reality and regulatory signaling, but intraday price and sentiment dynamics would enable more precise temporal decomposition.

Additional research avenues include:

\begin{itemize}
\item \textbf{High-frequency analysis:} Intraday data could reveal the temporal dynamics of infrastructure vs regulatory shocks, distinguishing immediate liquidity impacts from slower information diffusion processes and confirming whether mechanical disruptions propagate faster than informational shocks.

\item \textbf{Cross-market contagion:} Investigation of spillover effects between centralized exchanges, DeFi protocols, and traditional financial markets during infrastructure failures would clarify whether these events create systemic risks while regulatory events remain asset-specific.

\item \textbf{Regulatory heterogeneity:} Comparative analysis across jurisdictions (US SEC vs EU MiCA vs Asia-Pacific regimes) could test whether regulatory clarity moderates volatility responses, potentially explaining cross-sectional heterogeneity through differential regulatory exposure.

\item \textbf{Regulatory transmission channel decomposition:} A critical avenue for ongoing investigation involves decomposing the regulatory event channel to understand why these events generate smaller volatility impacts (0.419\% vs 2.385\%). Our findings demonstrate regulatory events have measurable effects, but the transmission mechanism remains unclear. Comparative analysis with traditional equity markets examines whether cryptocurrency regulatory responses operate through classical microstructure channels (bid-ask spreads, order book depth, trading volume) or purely through sentiment dynamics. Ongoing evidence demonstrates that cryptocurrency markets may exhibit a unique regulatory transmission mechanism where sentiment accounts for the entire effect with zero contribution from microstructure variables—a pattern not observed in equity or bond markets where approximately 60\% of regulatory impact transmits through microstructure changes. If validated, this would fundamentally challenge the applicability of traditional market microstructure theory to digital asset regulation, suggesting regulatory policy designed for traditional markets may operate through entirely different mechanisms in cryptocurrency contexts.

\item \textbf{Machine learning approaches:} Deep learning models trained on event embeddings could potentially improve event type classification beyond manual categorization, while NLP techniques might enable real-time event severity assessment moving beyond binary classifications.

\item \textbf{Volatility persistence dynamics:} Investigating whether near-integrated volatility represents a permanent characteristic of cryptocurrency markets or temporary phenomenon as markets mature would inform long-term risk modeling and forecasting strategies.
\end{itemize}

\subsection{Final Remarks}

Cryptocurrency markets continue evolving rapidly, yet our findings establish fundamental characteristics that appear structural rather than transitory. The 5.7$\times$ infrastructure-regulatory volatility multiplier, robust across multiple statistical tests and validation frameworks, demonstrates these markets exhibit sophisticated information processing capabilities that distinguish mechanical disruptions from gradual information absorption. The extreme volatility persistence documented (parameters approaching unity) represents a distinct regime requiring theoretical explanation and practical accommodation, fundamentally altering optimal risk management and forecasting strategies.

The superiority of asymmetric models with exogenous variables confirms that cryptocurrency volatility exhibits complex dynamics requiring sophisticated modeling approaches. TARCH-X specifications achieve superior information-theoretic fit for 83\% of assets, validating the inclusion of event-specific indicators and decomposed sentiment variables despite parsimony penalties. The methodological innovations -- custom TARCH-X maximum likelihood estimation, GDELT sentiment decomposition by event type proportions, and comprehensive multiple testing corrections -- provide a framework for future cryptocurrency event studies while demonstrating the feasibility of rigorous academic analysis in this rapidly developing domain.

The practical implications are substantial: portfolio managers allocating capital to cryptocurrency markets should employ differentiated hedging strategies for infrastructure versus regulatory risk, with infrastructure events requiring 4--5$\times$ higher capital buffers. The regime-switching findings amplify this urgency—during crisis periods, infrastructure sensitivity increases 5$\times$ (from 2.3\% to 11.2\%), requiring dynamic capital allocation that responds to market stress indicators rather than static risk multipliers. Traditional VaR frameworks assuming linear risk scaling will systematically underestimate tail risk during market drawdowns, precisely when adequate capitalization is most critical. Regulatory authorities should prioritize operational resilience standards given infrastructure failures generate larger market disruptions than policy announcements. Academic researchers examining cryptocurrency market dynamics should account for the unique near-integrated volatility regime and employ appropriate multiple testing corrections given the high event frequency in these markets.

As cryptocurrency markets mature toward greater institutional participation and regulatory integration, understanding their unique characteristics becomes increasingly critical. Our findings suggest the extreme persistence and infrastructure sensitivity may represent permanent structural features rather than temporary growing pains: fragmented exchange architecture, absence of designated market makers, and continuous 24/7 trading create conditions fundamentally different from traditional financial markets. Whether these characteristics persist or converge toward traditional market dynamics as institutions enter remains an open question with profound implications for market design, regulation, and global financial stability. Ongoing comparative analysis with equity markets investigates whether the observed regulatory response asymmetry reflects unique cryptocurrency market transmission mechanisms—if cryptocurrency regulatory events operate purely through sentiment channels while traditional markets exhibit substantial microstructure effects, this would establish digital assets as a fundamentally distinct asset class requiring entirely different regulatory frameworks.

This study provides empirical evidence and methodological tools for continued investigation of these essential questions at the intersection of technology and finance.

% ============================================================================
% REFERENCES (switch back to single-column)
% ============================================================================

\clearpage
\onecolumn

% ====================
% ACKNOWLEDGMENTS & DECLARATIONS
% ====================

\section*{Acknowledgments}

This research was conducted in accordance with the Charter of Farzulla Research, available at \url{https://farzulla.org}.

I am deeply grateful to my mother and sister for their unwavering support throughout this research journey. I acknowledge King's College London and the supervision provided during the initial stages of this research as part of my Master's dissertation, though this version significantly differs from and demonstrates substantial improvements over that earlier work.

The author acknowledges Perplexity AI for exceptional research discovery capabilities and Anthropic for Claude, whose invaluable assistance with analytical framework development, custom code implementations, methodological critique, and technical writing substantially accelerated this research and made it possible to expand the scope of analysis, implement manual TARCH-X estimation, and iteratively improve this work across multiple versions.

I am indebted to the open source community whose software tools (Python, NumPy, pandas, Matplotlib, and many others) made this research possible. Finally, I thank the many scholars whose work informed this framework - their insights into event study methodology, GARCH modeling, cryptocurrency markets, and financial econometrics provided the intellectual foundation upon which this analysis rests.

All errors, omissions, and interpretive limitations remain the author's responsibility.

\vspace{0.5em}
\noindent\textbf{Methodologies:} Research methodologies and reproducibility practices are documented at \href{https://farzulla.org/methodologies}{\texttt{farzulla.org/methodologies}}.

\section*{Funding}

This research received no external funding.

\section*{Conflict of Interest}

The author declares no conflict of interest.

\section*{Data Availability Statement}

All data, analysis code, figure generation scripts, and comprehensive documentation are publicly available at: \url{https://github.com/studiofarzulla/crypto-event-study}. The repository includes complete implementation details, parameter configurations, and instructions for reproducing all results. Data sources include CoinGecko historical price data and GDELT sentiment analysis. Models are implemented in Python 3.9+ and can be executed on standard hardware. The complete research package is archived on Zenodo with DOI: 10.5281/zenodo.17677682. An interactive dashboard for exploring all results is available at: \url{https://farzulla.org/research/crypto-event-study/}.

% ====================
% REFERENCES
% ====================

\bibliography{references}

% ====================
% APPENDICES
% ====================

\section*{Appendix A: Event List}
\addcontentsline{toc}{section}{Appendix A: Event List}

\textbf{Event Summary:} 50 events total (26 infrastructure, 24 regulatory) spanning January 2019 to August 2025, curated from an initial set of 207 events identified during research. The complete original event list is available in the GitHub repository. Events are listed chronologically below.

\vspace{1em}

\subsection*{2019}

\textbf{15 February:} QuadrigaCX exchange collapses after CEO death leaves private keys inaccessible (Infrastructure)

\textbf{3 April:} SEC publishes FinHub framework for digital asset classification (Regulatory)

\textbf{7 May:} Binance hack of 7,000 BTC, approximately USD 40 million (Infrastructure)

\textbf{18 June:} Facebook announces Libra stablecoin project (Regulatory)

\textbf{5 August:} Litecoin second halving (Infrastructure)

\textbf{24 October:} China President Xi Jinping endorses blockchain technology (Regulatory)

\subsection*{2020}

\textbf{12--13 March:} Black Thursday market crash triggers exchange outages (Infrastructure)

\textbf{11 May:} Third Bitcoin halving reduces block reward to 6.25 BTC (Infrastructure)

\textbf{15 June:} Compound token launch initiates DeFi summer (Infrastructure)

\textbf{1 September:} Binance Smart Chain mainnet launch (Infrastructure)

\textbf{1 December:} Ethereum 2.0 Beacon chain launch (Infrastructure)

\textbf{22 December:} SEC files lawsuit against Ripple Labs for XRP sales (Regulatory)

\subsection*{2021}

\textbf{8 February:} Tesla announces USD 1.5 billion Bitcoin purchase (Regulatory)

\textbf{14 April:} Coinbase direct listing on Nasdaq at USD 100 billion valuation (Infrastructure)

\textbf{19--21 May:} China announces cryptocurrency mining crackdown (Regulatory)

\textbf{9 June:} El Salvador adopts Bitcoin as legal tender (Regulatory)

\textbf{5 August:} Ethereum London hard fork implements EIP-1559 (Infrastructure)

\textbf{10 August:} Poly Network hack of USD 611 million (Infrastructure)

\textbf{24 September:} China announces total ban on cryptocurrency transactions (Regulatory)

\textbf{19 October:} ProShares Bitcoin Strategy ETF launches (Regulatory)

\textbf{14 November:} Bitcoin Taproot soft fork activates, enabling Schnorr signatures and MAST (Infrastructure)

\subsection*{2022}

\textbf{5--6 January:} Kazakhstan internet shutdown affects global mining (Infrastructure)

\textbf{9 March:} US President Biden issues executive order on digital assets (Regulatory)

\textbf{5--9 May:} Terra/Luna UST stablecoin collapse (Infrastructure)

\textbf{June:} Celsius Network and Three Arrows Capital failures (Infrastructure)

\textbf{15 September:} Ethereum Merge to proof-of-stake (Infrastructure)

\textbf{6 October:} BNB Chain bridge exploit of USD 570 million (Infrastructure)

\textbf{8--11 November:} FTX exchange bankruptcy and hack (Infrastructure)

\subsection*{2023}

\textbf{10--11 March:} Silicon Valley Bank collapse causes USDC depeg (Infrastructure)

\textbf{12 April:} Ethereum Shanghai upgrade enables staking withdrawals (Infrastructure)

\textbf{5--6 June:} SEC files lawsuits against Binance and Coinbase (Regulatory)

\textbf{15 June:} BlackRock files for spot Bitcoin ETF (Regulatory)

\textbf{29 August:} DC Circuit Court rules against SEC in Grayscale case (Regulatory)

\textbf{1 October:} European Union finalises MiCA regulation (Regulatory)

\textbf{21 November:} Binance settles with US authorities for USD 4.3 billion (Regulatory)

\subsection*{2024}

\textbf{10 January:} SEC approves eleven spot Bitcoin ETFs (Regulatory)

\textbf{13 March:} Ethereum Dencun upgrade implements proto-danksharding (Infrastructure)

\textbf{20 April:} Fourth Bitcoin halving reduces reward to 3.125 BTC (Infrastructure)

\textbf{23 May:} SEC approves spot Ethereum ETF rule changes (Regulatory)

\textbf{30 June:} EU MiCA Phase 1 implementation for stablecoins (Regulatory)

\textbf{23 July:} Spot Ethereum ETFs begin trading (Infrastructure)

\subsection*{2025}

\textbf{21 February:} Bybit exchange hack of USD 1.5 billion (Infrastructure)

\textbf{27 February:} SEC and Coinbase file joint stipulation to dismiss case (Regulatory)

\textbf{7 March:} OCC Interpretive Letter 1183 reaffirms permissible bank crypto activities (custody, certain stablecoin operations) and rescinds prior limits (Regulatory)

\textbf{4 April:} SEC clarifies stablecoins not securities (Regulatory)

\textbf{7 May:} Ethereum Prague-Electra upgrade (Infrastructure)

\textbf{18 July:} GENIUS Act enacted as first US federal stablecoin framework (Regulatory)

\textbf{29 July:} SEC allows in-kind creations/redemptions for crypto ETPs (Regulatory)

\textbf{8 August:} SEC v Ripple litigation concludes favouring Ripple (Regulatory)

\section*{Appendix B: GDELT Data Extraction}
\addcontentsline{toc}{section}{Appendix B: GDELT Data Extraction}

The GDELT sentiment data was extracted using a multi-stage SQL query in Google BigQuery, processing the Global Knowledge Graph database from January 2019 to August 2025. The complete query implements three-stage processing:

\begin{enumerate}[leftmargin=*,nosep]
\item \textbf{Filtering:} Cryptocurrency-related articles identified through theme keywords (bitcoin, crypto, ethereum)
\item \textbf{Sentiment calculation:} Volume-weighted sentiment scores using article tone and mention counts
\item \textbf{Normalization:} 52-week rolling z-score standardization with theme decomposition (regulatory vs infrastructure)
\end{enumerate}

The full SQL implementation is available in the project repository.

\section*{Appendix C: TARCH-X Implementation}
\addcontentsline{toc}{section}{Appendix C: TARCH-X Implementation}

Given limitations in existing econometric software for implementing exogenous variables directly in the variance equation, this study developed a custom maximum likelihood estimator. The implementation ensures precise specification of the theoretical model where conditional variance follows:

\begin{align*}
\sigma^2_t &= \omega + \alpha \cdot \varepsilon^2_{t-1} + \gamma \cdot \varepsilon^2_{t-1} \cdot I(\varepsilon_{t-1} < 0) \\
&\quad + \beta \cdot \sigma^2_{t-1} + \sum_j \delta_j \cdot x_{j,t}
\end{align*}

The manual implementation provides full control over the optimisation process, transparent likelihood function specification, and proper computation of robust standard errors via numerical Hessian. The complete Python implementation spans approximately 400 lines and includes parameter constraint handling, Student-$t$ likelihood computation, and bootstrap inference capabilities.


\section*{Appendix D: Hypothesis Outcome Changes}
\addcontentsline{toc}{section}{Appendix D: Hypothesis Outcome Changes}

\begin{table*}[htbp]
\centering
\caption{Hypothesis Outcome Changes}
\label{tab:hypothesis_changes}
\footnotesize
\begin{tabular}{@{}lp{2.2cm}p{2.5cm}@{}}
\toprule
\textbf{Hypothesis} & \textbf{OLD Status} & \textbf{NEW Status} \\
\midrule
H1: Infra $>$ Reg & Rejected ($p=0.997$) & \textbf{Supported} ($p=0.0008$) \\
H2: Sentiment leading & Rejected & \textbf{Partial Support} \\
H3: TARCH-X superiority & Mixed & \textbf{Supported} (83\% AIC) \\
\bottomrule
\end{tabular}
\end{table*}

\clearpage
\onecolumn
\section*{Appendix E: Robustness Checks}
\addcontentsline{toc}{section}{Appendix E: Robustness Checks}

The 5.7$\times$ infrastructure-regulatory multiplier is robust across multiple validation approaches.

\begin{figure}[htbp]
\centering
\includegraphics[width=\textwidth]{robustness_effects_confidence_intervals.png}
\caption{Infrastructure vs Regulatory Effects with 95\% Confidence Intervals}
\label{fig:robustness_ci}

\footnotesize
Error bars represent 95\% confidence intervals calculated from parameter standard errors. Infrastructure effects (blue) are significantly larger and more precisely estimated than regulatory effects (purple) across all six cryptocurrencies. Non-overlapping confidence intervals provide additional evidence for the infrastructure-regulatory differential beyond point estimates.
\end{figure}


\begin{figure}[htbp]
\centering
\includegraphics[width=\textwidth]{placebo_test_robustness.png}
\caption{Placebo Test Results}
\label{fig:placebo}

\footnotesize
Distribution of placebo heterogeneity statistics from 1,000 randomly permuted event samples. The observed infrastructure-regulatory differential (red vertical line at 1.97pp) significantly exceeds the null distribution, confirming that the effect is driven by genuine event-type differences rather than random variation. The p-value from the placebo test is p < 0.001, providing additional confirmation of statistical significance.
\end{figure}

\end{document}

% ============================================================================
% VERSION CONTROL GUIDE
% ============================================================================
%
% Semantic Versioning for Academic Papers:
%   vMAJOR.MINOR.PATCH
%
% MAJOR: Complete rewrite, new methodology, fundamentally different paper
%        Example: 1.0.0 -> 2.0.0 (changed from single case study to multi-country)
%
% MINOR: New experiments, expanded scope, additional sections
%        Example: 1.0.0 -> 1.1.0 (added robustness checks, new case study)
%
% PATCH: Bug fixes, typo corrections, citation updates, clarifications
%        Example: 1.0.0 -> 1.0.1 (fixed table formatting, corrected citations)
%
% To update version:
%   1. Change \newcommand{\paperver}{X.Y.Z} above
%   2. Update \newcommand{\paperdate}{Month Year}
%   3. Update "Note on Prior Work" section describing changes
%   4. Recompile PDF
%   5. Upload new version to Zenodo (creates linked DOI)
%
% ============================================================================
